% ============================================================
% main.tex — Memoria del Trabajo de Integración Curricular
%   "Modelado experimental de las fases prefill y decode de
%    Qwen3-8B en NVIDIA A100 mediante vLLM"
%   Código institucional: ITI-TFG-B26-0017
%   Autor: José Andrés Ureta Granda
%   Tutor: Andrés Xavier Rubio Proaño
%   Universidad Tecnológica Indoamérica — Facultad de Ingenierías
%   Ingeniería en Tecnologías de la Información
%
% Requiere en la carpeta del proyecto:
%   - uti_tfm.sty        (formato institucional UTI)
%   - referencias.bib    (vLLM/PagedAttention, Orca, Sarathi-Serve,
%                         DistServe, Splitwise, LLM-Viewer, Qwen3, ...)
%   - Logos/LOGO-UTI-1.png
%   - SPEC.md            (especificación del proyecto; fuente de los
%                         contenidos de cada capítulo, ver su sección 15)
%
% Compilación recomendada (por Century Gothic):
%   XeLaTeX → BibTeX → XeLaTeX → XeLaTeX
%   (requiere la fuente Century Gothic instalada en el sistema)
% Alternativa sin Century Gothic: latexmk -pdf main.tex
%   (pdfLaTeX; el .sty usa TeX Gyre Adventor como sustituto)
% ============================================================

\documentclass[12pt,a4paper,oneside]{report}

% Estilo UTI
\usepackage{uti_tfm}

% Referencias IEEE desde referencias.bib (BibTeX + IEEEtran)
\usepackage{cite}                        % ordena y comprime citas IEEE
\renewcommand{\citepunct}{], [}         % [1], [2]
\renewcommand{\citedash}{]--[}          % [1]--[3]
\usepackage[hidelinks]{hyperref}
% Compatibilidad con FirmaEC: evita que hyperref inserte /OpenAction.
\hypersetup{pdfstartview=}
\renewcommand{\bibname}{REFERENCIAS} % clase report

\begin{document}

% ==========================================================
% PORTADA (sin numeración visible; NO cuenta para la numeración)
% ==========================================================
% Título en inglés utilizado en la página de Abstract (traducción del
% título de la portada; validar con el tutor, SPEC §16.6).
\UTIDatosTituloIngles{Experimental Modeling of the Prefill and Decode Phases of Qwen3-8B on NVIDIA A100 Using vLLM}
% Lectores: nombre y género (masculino | femenino) de cada lector/a.
% PENDIENTE: completar cuando la Facultad designe a los lectores.
\UTIDatosLector1{Nombres Apellidos}{masculino}
\UTIDatosLector2{Nombres Apellidos}{femenino}
%
% LÍNEA Y SUBLÍNEA DE INVESTIGACIÓN
% Se selecciona con: \UTISeleccionInvestigacion{línea}{sublínea}
% Estado: seleccionada la propuesta; falta la confirmación del tutor
% (SPEC §16.6).
%   Seleccionada: {4}{6} Optimización del rendimiento
%                 (modelado de latencia y throughput de prefill/decode).
%   Alternativa:  {1}{3} Sistemas de almacenamiento y computación en la nube
%                 (dimensionamiento de un servicio de LLM sobre GPU).
%
% Catálogo completo:
%
% LÍNEA 1: Análisis, diseño e implementación de infraestructura tecnológica
%   {1}{1} Diseño y optimización de redes de comunicación.
%   {1}{2} Ciberseguridad y protección de infraestructuras tecnológicas.
%   {1}{3} Sistemas de almacenamiento y computación en la nube.  <- alternativa
%   {1}{4} Automatización y gestión de infraestructuras inteligentes.
%
% LÍNEA 2: Redes informáticas de comunicación
%   {2}{1} Arquitecturas de redes avanzadas.
%   {2}{2} Protocolos y estándares de comunicación.
%   {2}{3} Ciberseguridad en redes.
%   {2}{4} Redes inalámbricas e Internet de las Cosas (IoT).
%
% LÍNEA 3: Gestión de la información
%   {3}{1} Almacenamiento y recuperación de información.
%   {3}{2} Análisis y procesamiento de datos.
%   {3}{3} Gobernanza y seguridad de la información.
%   {3}{4} Ciclo de vida de la información.
%   {3}{5} Sistemas de gestión del conocimiento.
%   {3}{6} Visualización de información.
%   {3}{7} Tecnologías emergentes en la gestión de información.
%   {3}{8} Gestión de datos no estructurados.
%
% LÍNEA 4: Diseño, desarrollo e integración de sistemas basados en
%          arquitecturas web, móviles y sistemas embebidos
%   {4}{1} Diseño y desarrollo de aplicaciones web.
%   {4}{2} Desarrollo de aplicaciones móviles.
%   {4}{3} Integración de sistemas embebidos.
%   {4}{4} Interoperabilidad entre plataformas.
%   {4}{5} Seguridad en sistemas web, móviles y embebidos.
%   {4}{6} Optimización del rendimiento.                         <- seleccionada
%   {4}{7} Aplicación de tecnologías emergentes.
%   {4}{8} Sistemas basados en la nube y edge computing.
%
% Si el tutor elige la alternativa, cambiar a
% \UTISeleccionInvestigacion{1}{3}. El valor se imprime en la sección
% "Área de Estudio" del Capítulo II.
\UTISeleccionInvestigacion{4}{6}
\UTIPortada
{LOGO-UTI-1.png} % Logo (se busca en Logos/)
{FACULTAD DE INGENIERÍAS} % Facultad
{INGENIERÍA EN TECNOLOGÍAS DE LA INFORMACIÓN} % Carrera
{Modelado experimental de las fases prefill y decode de Qwen3-8B en NVIDIA A100 mediante vLLM} % Tema / Título
{Trabajo de Titulación modalidad Trabajo de Integración Curricular} % Modalidad
% Autor único del proyecto. \UTIDatosAutores recibe: nombre, género,
% cédula, dirección, correo electrónico y teléfono.
{\UTIDatosAutores
  {José Andrés Ureta Granda}{masculino}{1751545185}
  {Pichincha, Quito, Belisario Quevedo, La Gasca}{andredks334@gmail.com}{0939116995}} % Autor
{\UTIDatosTutor{Andrés Xavier Rubio Proaño}{masculino}} % Tutor
{Quito, Ecuador} % Ciudad, País
{2026} % Año

% ==========================================================
% PÁGINAS PRELIMINARES (romanos minúscula: i, ii, iii...)
% ==========================================================
\UTIFrontMatter

% Autorización para el Repositorio Digital Institucional (RDI-UTI)
\UTIFrontTitle{Autorización por parte del autor para la consulta, reproducción parcial o total, y publicación electrónica del Trabajo de Integración Curricular}

\UTIAutoresPronombre, \UTIAutoresTexto, \UTIVerboDeclaro{} ser \UTIAutorEtiquetaMin{} del Trabajo de Titulación con el nombre ``\UTITituloProyecto'', como requisito para optar \UTIFraseOptarTituloProfesional{} y \UTIVerboAutorizo{} al Sistema de Bibliotecas de la Universidad Tecnológica Indoamérica para que, con fines netamente académicos, divulgue esta obra a través del Repositorio Digital Institucional (RDI-UTI).

Los usuarios del RDI-UTI podrán consultar el contenido de este trabajo en las redes de información del país y del exterior, con las cuales la Universidad tenga convenios. La Universidad Tecnológica Indoamérica no se hace responsable por el plagio o copia del contenido parcial o total de este trabajo.
	        
Del mismo modo, \UTIVerboAcepto{} que los Derechos de Autor, Morales y Patrimoniales sobre esta obra serán compartidos entre \UTIAutoresPoseedor{} y la Universidad Tecnológica Indoamérica, y que no \UTIVerboTramito{} la publicación de este trabajo en ningún otro medio sin autorización expresa de la misma. En caso de que exista el potencial de generación de beneficios económicos o patentes, producto de este trabajo, \UTIVerboAcepto{} que se deberán firmar convenios específicos adicionales donde se acuerden los términos de adjudicación de dichos beneficios.

Para constancia de esta autorización, en \UTICiudadPais, a los xx días del mes de xx de \UTIAnio, \UTIVerboFirmo{} conforme:

\vspace{1cm}

\UTIRepositorioAutoresBloque

% Aprobación del tutor
\UTIIfTutorFemaleTF{\UTIFrontTitle{Aprobación de la tutora}}{\UTIFrontTitle{Aprobación del tutor}}

En mi calidad de \UTITutorEtiqueta{} del Trabajo de Titulación ``\UTITituloProyectoMayusculas'', presentado por \UTIAutoresTexto, para optar por \UTIFraseTituloProfesional{}.

\begin{center}
CERTIFICO
\end{center}

Que dicho Trabajo de Titulación ha sido revisado en todas sus partes y considero que reúne los requisitos y méritos suficientes para ser sometido a la presentación pública y evaluación por parte del Tribunal revisor que se designe.

\vspace{1.5cm}

{\raggedleft
\UTICiudadPais, xx de xx de \UTIAnio\par}

\vspace{2.5cm}

\UTIFirmaTutorBloque


% Declaración de autenticidad del autor
\UTIFrontTitle{Declaración de autenticidad}

\UTIAutoresPronombre, \UTIAutoresTexto, \UTIVerboDeclaro{} que los contenidos y los resultados obtenidos en el presente Trabajo de Titulación ``\UTITituloProyecto'', como requerimiento previo para \UTIFraseObtencionTituloProfesional{}, son absolutamente originales, auténticos y personales y de exclusiva responsabilidad legal y académica \UTIAutoresResponsabilidad.

\vspace{1.5cm}

{\raggedleft
\UTICiudadPais, xx de xx de \UTIAnio\par}

\vspace{2.5cm}

\UTIFirmasAutoresBloque


% Aprobación de lectores (depende de \UTIDatosLector1/2)
\UTIFrontTitle{Aprobación de lectores}

El Trabajo de Titulación ha sido revisado, aprobado y autorizada su impresión y empastado sobre el tema ``\UTITituloProyectoMayusculas'', presentado por \UTIAutoresTexto, previo a \UTIFraseObtencionTituloProfesional{}, y reúne los requisitos de fondo y forma para su sustentación.

\vspace{1.5cm}
\UTICiudadPais, xx de xx de \UTIAnio
\vspace{2.5cm}

\UTILectoresFirmaBloque

\UTIIfTwoAuthorsTF{%
\UTIFrontTitle{Dedicatoria}

\vfill

\begin{flushright}
\textbf{\UTIAutorUnoNombre}\\
(Obligatoria)\\
La dedicatoria de este autor o autora va en la parte inferior derecha de la página.
\end{flushright}

\UTIFrontTitle{Dedicatoria}

\vfill

\begin{flushright}
\textbf{\UTIAutorDosNombre}\\
(Obligatoria)\\
La dedicatoria de este autor o autora va en la parte inferior derecha de la página.
\end{flushright}
}{%
\UTIFrontTitle{Dedicatoria}

\vfill

\begin{flushright}
(Obligatoria)\\
La dedicatoria va en la parte inferior derecha de la página.
\end{flushright}
}

\UTIIfTwoAuthorsTF{%
\UTIFrontTitle{Agradecimiento}

\vfill

\begin{flushright}
\textbf{\UTIAutorUnoNombre}\\
(Obligatorio)\\
El agradecimiento de este autor o autora va en la parte inferior derecha de la página.
\end{flushright}

\UTIFrontTitle{Agradecimiento}

\vfill

\begin{flushright}
\textbf{\UTIAutorDosNombre}\\
(Obligatorio)\\
El agradecimiento de este autor o autora va en la parte inferior derecha de la página.
\end{flushright}
}{%
\UTIFrontTitle{Agradecimiento}

\vfill

\begin{flushright}
(Obligatorio)\\
El agradecimiento va en la parte inferior derecha de la página.
\end{flushright}
}

% --------------------------
% ÍNDICES (contenidos, tablas, figuras, etc.; según el .sty)
% --------------------------
% Índice de contenidos
\clearpage
\tableofcontents
% Índice de tablas
\clearpage
\listoftables
% Índice de figuras
\clearpage
\listoffigures
% Índice de imágenes
\clearpage
\listofimagenes
% Índice de ecuaciones
\clearpage
\listofecuaciones
% Índice de anexos
\clearpage
\listofanexos
\clearpage

% Resumen: fases prefill/decode, banco de pruebas con vLLM en A100,
% modelos predictivos de TTFT, TPOT y throughput, y su validación.
\UTIResumenPagina
  {spanish}
  {Resumen Ejecutivo}
  {\textit{[Redacte el resumen del Trabajo de Integración Curricular en un solo párrafo, con un máximo de 300 palabras. Exponga de manera clara el problema, el objetivo general, la metodología empleada, los principales resultados obtenidos y las conclusiones.]}}
  {Descriptores}
  {\textit{[Incluya de 3 a 4 palabras clave representativas, ordenadas alfabéticamente y separadas por comas.]}}

% Abstract (usa el título en inglés definido arriba)
\UTIResumenPagina
  {english}
  {Abstract}
  {\textit{[Write the English version of the executive summary in a single paragraph, with a maximum of 300 words. Clearly present the problem, general objective, methodology, main results, and conclusions.]}}
  {Keywords}
  {\textit{[Include 3 to 4 representative keywords in alphabetical order, separated by commas.]}}

% ==========================================================
% CUERPO (arábigos CONTINUANDO el contador desde los romanos)
% Nota: el .sty cambia el formato de página a arábigos sin reiniciar.
% ==========================================================
\UTIMainMatter

% ==========================================================
% CAPÍTULO I — INTRODUCCIÓN
%   Contextualización, problema y justificación: dimensionamiento de
%     servicios de LLM y métricas agregadas que mezclan prefill y decode
%     (SPEC §1–2).
%   Estado del arte (>= 5 trabajos IEEE): vLLM/PagedAttention, Orca,
%     Sarathi-Serve, DistServe, Splitwise, LLM-Viewer/roofline, Qwen3.
%   Objetivos OE1–OE4 y alcance: un modelo (Qwen3-8B, BF16), un motor
%     (vLLM) y una GPU NVIDIA A100 (SPEC §3–4).
%   Fundamentación teórica: prefill/decode, caché KV, continuous batching,
%     chunked prefill, modelo roofline, MFU y MBU (SPEC §6).
% ==========================================================
\chapter{Introducción}
\label{chap:introduccion}

\section{Antecedentes}
\label{sec:cap1:antecedentes}

\subsection{Contexto}
\label{sec:cap1:contexto}

Los modelos de lenguaje de gran tamaño (LLM, por sus siglas en inglés) se han incorporado a asistentes conversacionales, sistemas de consulta documental y flujos de procesamiento automático de texto. La mayoría de estos modelos se basa en la arquitectura Transformer~\cite{vaswani2017attention}. Una vez entrenado, cada modelo se ejecuta en modo de inferencia ante cada solicitud recibida. El costo de servirlo se acumula, por tanto, durante toda su vida operativa.

La inferencia de un LLM se desarrolla en dos fases con comportamientos distintos. La fase prefill procesa el texto de entrada de la solicitud, o prompt, y produce el primer token de salida. La fase decode genera los tokens restantes de uno en uno, de forma autorregresiva. Durante ambas fases, el modelo conserva en la memoria de la unidad de procesamiento gráfico (GPU) las claves y los valores de atención de los tokens previos; esta estructura se denomina KV cache.

La distinción entre prefill y decode ocupa un lugar central en los trabajos revisados en la Sección~\ref{sec:cap1:estado-arte}. Cada fase somete a la GPU a exigencias diferentes. La configuración de inferencia elegida determina el reparto de recursos entre ambas fases y condiciona la experiencia de los usuarios.

El presente trabajo de fin de grado (TFG) se enmarca en el proyecto LLM-CoDesign. Este proyecto caracteriza la inferencia de modelos de lenguaje considerando conjuntamente el modelo, la carga de trabajo, el hardware y el software. Su alcance general comprende múltiples modelos, GPU, motores y métricas durante cinco años. El TFG reduce ese alcance a un caso experimental controlado y constituye el primer caso de estudio del proyecto.

El caso seleccionado combina el modelo Qwen3-8B, perteneciente a la familia Qwen3~\cite{qwen2025qwen3}, con una GPU NVIDIA A100 y el motor de inferencia vLLM. La A100 se comercializa en variantes de 40~GB y 80~GB de memoria~\cite{nvidia2021a100}.

\subsection{Estado del arte}
\label{sec:cap1:estado-arte}

La literatura sobre inferencia de LLM se concentra en mecanismos de sistema orientados a aprovechar mejor la GPU. Las líneas principales abarcan la planificación de solicitudes, la gestión de memoria, los kernels de atención, el tratamiento diferenciado de las fases y la cuantización. Esta sección revisa seis trabajos representativos de esas líneas y los relaciona con el presente TFG.

Yu et al.~\cite{yu2022orca} proponen Orca, un motor de inferencia que planifica la ejecución a nivel de iteración en lugar de a nivel de solicitud. Con esta estrategia, el sistema incorpora solicitudes nuevas al lote y retira las finalizadas en cada paso de generación. Orca reserva el espacio del KV cache de cada solicitud en bloques contiguos dimensionados de antemano; esta reserva desperdicia una parte de la memoria disponible~\cite{kwon2023pagedattention}.

Kwon et al.~\cite{kwon2023pagedattention} presentan PagedAttention, un mecanismo que gestiona el KV cache en bloques de tamaño fijo, de forma análoga a la memoria virtual paginada de un sistema operativo. La gestión paginada reduce la fragmentación de la memoria y permite compartir bloques entre solicitudes. Sobre este mecanismo, los autores construyen vLLM, un motor de inferencia que combina la gestión paginada con la planificación a nivel de iteración, conocida como batching continuo. vLLM es el motor de inferencia empleado en este trabajo.

Dao et al.~\cite{dao2022flashattention} desarrollan FlashAttention, un algoritmo de atención exacta que reduce los accesos a la memoria de alto ancho de banda de la GPU mediante el procesamiento por bloques. Su aporte pone de relieve el peso del movimiento de datos entre niveles de memoria en el costo de la atención. Esta perspectiva resulta pertinente para interpretar los cuellos de botella de cada fase.

Agrawal et al.~\cite{agrawal2024sarathi} estudian el compromiso entre throughput y latencia en la inferencia de LLM; el throughput se entiende aquí como el número de tokens procesados por segundo. Los autores observan que la incorporación del prefill de solicitudes nuevas a un lote interrumpe la generación de las solicitudes en decode. Los autores proponen Sarathi-Serve como respuesta; este motor de inferencia divide el prefill en fragmentos (chunked prefill) y los planifica sin detener la generación en curso. La coexistencia de ambas fases en un mismo lote condiciona la manera de medirlas por separado.

Zhong et al.~\cite{zhong2024distserve} proponen DistServe, un motor de inferencia que asigna prefill y decode a GPU distintas para eliminar la interferencia entre fases. Su análisis caracteriza prefill como una fase limitada por el cómputo y decode como una fase limitada por el ancho de banda de memoria. Los autores evalúan cada fase con una métrica propia: el TTFT (\textit{time to first token}), tiempo hasta el primer token, asociado a prefill; y el TPOT (\textit{time per output token}), tiempo por token de salida, asociado a decode.

Lin et al.~\cite{lin2024awq} proponen AWQ, un método de cuantización de solo pesos a baja precisión. El método protege los pesos más relevantes a partir de la distribución de las activaciones. La representación de los pesos con 4 bits (AWQ INT4) reduce la memoria ocupada por el modelo frente al formato de coma flotante de 16 bits BF16. Este TFG compara ambas representaciones numéricas.

La Tabla~\ref{tab:cap1:estado-arte} resume el aporte, la limitación principal y la relación de cada trabajo con el presente TFG. Los seis trabajos proponen o evalúan mecanismos de sistema; ninguno caracteriza la combinación de Qwen3-8B, A100 y vLLM. Todos ellos son anteriores a la publicación de la familia Qwen3~\cite{qwen2025qwen3}.

\begin{table}[htbp]
  \centering
  \caption{Comparación de los trabajos revisados en el estado del arte.}
  \label{tab:cap1:estado-arte}
  \small
  \begin{tabular}{|p{0.17\textwidth}|p{0.23\textwidth}|p{0.23\textwidth}|p{0.23\textwidth}|}
    \hline
    \textbf{Trabajo} & \textbf{Aporte principal} & \textbf{Limitación principal} & \textbf{Relación con el TFG} \\
    \hline
    Orca~\cite{yu2022orca} & Planificación a nivel de iteración & Reserva contigua del KV cache & Origen del batching continuo de vLLM \\
    \hline
    vLLM~\cite{kwon2023pagedattention} & Gestión paginada del KV cache & No separa el análisis por fases & Motor de inferencia del TFG \\
    \hline
    FlashAttention \cite{dao2022flashattention} & Atención exacta con menos accesos a memoria & Optimiza un kernel aislado & Interpretación de cuellos de botella \\
    \hline
    Sarathi-Serve~\cite{agrawal2024sarathi} & Chunked prefill sin detener decode & Mezcla ambas fases en un mismo lote & Condiciona la medición por fases \\
    \hline
    DistServe~\cite{zhong2024distserve} & Separación de prefill y decode en GPU distintas & Requiere varias GPU & Define TTFT y TPOT por fase \\
    \hline
    AWQ~\cite{lin2024awq} & Cuantización de pesos a 4 bits & Posible degradación de calidad & Representación AWQ INT4 evaluada \\
    \hline
  \end{tabular}
\end{table}
\elaboradopor{José Andrés Ureta Granda}
\fuente{Elaboración propia}

\section{Planteamiento del problema}
\label{sec:cap1:planteamiento}

La selección de una configuración de inferencia suele basarse en métricas agregadas, como los tokens por segundo, sin distinguir el comportamiento de las fases prefill y decode. Una métrica agregada resume en una sola cifra dos procesos con exigencias distintas sobre la GPU~\cite{zhong2024distserve}. El resultado es una visión incompleta del sistema.

Esta práctica impide explicar varios fenómenos relevantes para el despliegue. No permite justificar por qué una configuración mejora el throughput y empeora la latencia interactiva. La práctica deja sin aclarar en qué casos la cuantización permite aumentar la concurrencia, entendida como el número de solicitudes simultáneas. Deja sin explicar, además, la influencia de la longitud de entrada, es decir, el número de tokens del prompt, sobre el TTFT; y la del crecimiento del KV cache sobre la latencia durante decode.

La consecuencia práctica es la dificultad para elegir una configuración adecuada según el uso previsto. Un servicio de chat interactivo, una consulta documental basada en recuperación (RAG) y un procesamiento por lotes imponen prioridades diferentes. Una configuración apropiada para un escenario puede resultar inadecuada para otro.

Qwen3-8B puede ejecutarse sin dificultad en una A100. Esta posibilidad no garantiza la existencia de una configuración universalmente óptima. El rendimiento depende de la carga de trabajo, del nivel de concurrencia y de la representación numérica del modelo; para esta combinación concreta no se dispone de una caracterización experimental que separe ambas fases.

\section{Formulación del problema}
\label{sec:cap1:formulacion}

El problema se formula mediante la siguiente pregunta principal: ¿cómo influyen la cuantización, la longitud de contexto y la concurrencia en el comportamiento de las fases prefill y decode de Qwen3-8B, ejecutado mediante vLLM sobre una GPU NVIDIA A100, en términos de rendimiento, memoria, energía y calidad?

La longitud de contexto se controla mediante la longitud de entrada y la longitud de salida, definida como el número de tokens generados por solicitud. De la pregunta principal se derivan las siguientes preguntas específicas:

\begin{enumerate}
  \item ¿Cómo cambia el TTFT cuando aumenta la longitud de entrada?
  \item ¿Cómo cambia el TPOT cuando aumenta la concurrencia?
  \item ¿Cómo influye la longitud de salida en la latencia y en la energía?
  \item ¿Cuánta memoria ocupan las representaciones numéricas BF16 y AWQ INT4?
  \item ¿En qué condiciones la cuantización mejora la capacidad del sistema?
  \item ¿Qué configuraciones producen saturación o errores por falta de memoria?
  \item ¿Cuáles son los principales cuellos de botella de cada fase?
  \item ¿Con qué error estima un modelo predictivo básico el TTFT y el TPOT a partir de las variables experimentales?
\end{enumerate}

\section{Objetivos}
\label{sec:cap1:objetivos}

\subsection{Objetivo general}
\label{sec:cap1:objetivo-general}

Caracterizar y modelar experimentalmente el comportamiento de las fases prefill y decode de Qwen3-8B sobre una GPU NVIDIA A100, evaluando el efecto de la cuantización, la longitud de contexto y la concurrencia sobre el rendimiento, la memoria, la energía y la calidad.

\subsection{Objetivos específicos}
\label{sec:cap1:objetivos-especificos}

{\setlength{\emergencystretch}{1em}% evita desbordes en los OE literales
\begin{itemize}
  \item[OE1.] Implementar un entorno reproducible para ejecutar Qwen3-8B mediante vLLM sobre una A100.
  \item[OE2.] Diseñar un benchmark parametrizable para generar cargas con diferentes longitudes de entrada, salida y niveles de concurrencia.
  \item[OE3.] Medir TTFT, TPOT, throughput, memoria GPU, utilización y consumo energético.
  \item[OE4.] Analizar separadamente las fases prefill y decode e identificar sus principales cuellos de botella.
  \item[OE5.] Construir un modelo predictivo básico para estimar TTFT y TPOT.
\end{itemize}}

Para OE5, la ficha institucional establece como meta referencial un error porcentual absoluto medio (MAPE) inferior al 15--20~\% en las métricas principales. Esta meta es ajustable después de los experimentos piloto. No constituye, por tanto, un compromiso cerrado sobre la precisión del modelo predictivo.

\section{Justificación}
\label{sec:cap1:justificacion}

Desde el punto de vista técnico, el análisis separado de prefill y decode aporta una explicación de los compromisos entre latencia, throughput y memoria. Las métricas agregadas registran estos compromisos sin explicarlos. Williams et al.~\cite{williams2009roofline} proponen el modelo roofline, que relaciona el rendimiento alcanzable con los límites de cómputo y de ancho de banda de memoria del procesador. Este marco ofrece una referencia para situar cada fase respecto a los límites de la A100 y sustenta el análisis de cuellos de botella de OE4.

Desde el punto de vista operativo, el trabajo orienta la elección de configuraciones para tres escenarios de uso. El chat interactivo prioriza un TTFT y un TPOT bajos. La consulta documental requiere un equilibrio entre longitud de entrada, latencia y memoria; el procesamiento por lotes prioriza el throughput y la eficiencia energética.

La utilidad del trabajo se materializa en el prototipo Qwen-A100 Profiler. Esta herramienta automatiza la ejecución de los experimentos, recolecta las métricas, almacena los resultados de forma estructurada y estima el TTFT y el TPOT mediante el modelo predictivo de OE5. El conjunto de datos experimental y los procedimientos de reproducción completan los entregables reutilizables.

El beneficiario directo es el proyecto LLM-CoDesign, que obtiene su primer caso de estudio y una metodología de caracterización aplicable a otras combinaciones de modelo, hardware y software. Los beneficiarios indirectos son los equipos técnicos que despliegan LLM en infraestructura propia y necesitan criterios cuantitativos para configurar el servicio.

El trabajo es factible en el plano técnico. Qwen3-8B cabe en la memoria de una única A100, y las herramientas necesarias (vLLM, NVML, DCGM y las bibliotecas de análisis de datos) son de libre acceso. El trabajo no requiere entrenar el modelo, desarrollar kernels de GPU ni modificar internamente vLLM. En el plano operativo, el plan de trabajo comprende doce semanas con una dedicación de 15 a 20 horas semanales; la factibilidad depende de disponer de acceso controlado a la GPU durante las campañas de medición.

\section{Alcance y delimitaciones}
\label{sec:cap1:alcance}

\subsection{Alcance}
\label{sec:cap1:alcance-incluye}

La unidad de análisis es la configuración experimental $C = (M, W, H, S)$. En ella, $M$ es el modelo Qwen3-8B; $W$, las características de la carga de trabajo; $H$, una GPU NVIDIA A100; y $S$, vLLM con sus parámetros de configuración. La variante exacta de la A100, de 40~GB u 80~GB, se documenta como parte del entorno experimental, porque su capacidad de memoria condiciona los niveles máximos de contexto y concurrencia alcanzables.

Las variables experimentales y sus niveles son los siguientes:

\begin{itemize}
  \item Representación numérica: BF16 y AWQ INT4 (2 niveles).
  \item Longitud de entrada: 128, 512, 2048 y 8192 tokens (4 niveles).
  \item Longitud de salida: 64 y 256 tokens (2 niveles).
  \item Concurrencia: 1, 4, 8 y 16 solicitudes simultáneas (4 niveles).
\end{itemize}

La combinación de los niveles produce $2 \times 4 \times 2 \times 4 = 64$ configuraciones. Con cinco repeticiones por configuración, la campaña comprende alrededor de 320 ejecuciones medidas. Las configuraciones se agrupan además en tres tipos de carga de trabajo: chat corto, consulta documental y contexto largo.

Las métricas registradas son el TTFT, el TPOT, el throughput, la latencia total de la solicitud, la memoria máxima utilizada, la utilización de la GPU, la potencia media y la energía estimada por token. La energía se presenta como una estimación basada en la telemetría interna de la GPU, obtenida mediante NVML y DCGM. No constituye una medición eléctrica calibrada externamente.

El alcance incluye una comparación limitada de calidad entre BF16 y AWQ INT4 para identificar posibles degradaciones asociadas a la cuantización. Incluye también el prototipo Qwen-A100 Profiler en su versión de producto mínimo viable, con una recomendación básica de configuraciones para los tres escenarios de uso.

\subsection{Delimitaciones}
\label{sec:cap1:delimitaciones}

El trabajo excluye de forma explícita los siguientes elementos:

\begin{itemize}
  \item Entrenamiento completo o fine-tuning del modelo.
  \item Comparación con otras familias de LLM.
  \item Evaluación de múltiples GPU y ejecución multi-GPU.
  \item Paralelismo tensorial o de pipeline.
  \item Comparación entre varios motores de inferencia.
  \item Modelos Mixture-of-Experts o multimodales.
  \item Speculative decoding.
  \item Longitudes de entrada superiores a 8192 tokens en la campaña principal.
  \item Predicción del comportamiento en otro hardware.
  \item Simulación de sistemas distribuidos.
  \item Desarrollo completo del asesor multiobjetivo de LLM-CoDesign.
\end{itemize}

Las conclusiones del trabajo serán válidas únicamente para la combinación de Qwen3-8B, NVIDIA A100 y vLLM. No se formularán afirmaciones de superioridad general sobre otras plataformas o modelos.

% PENDIENTE (opción (a) aprobada por el autor): esta sección se redacta en
% una sesión posterior. Aquí irá el marco teórico (Transformer, GQA
% [ainslie2023gqa], KV cache, PagedAttention, AWQ, roofline, etc.).
\section{Fundamentación teórica}
\label{sec:cap1:fundamentacion}
\textit{[Desarrollar los fundamentos teóricos, conceptos, modelos y tecnologías que sustentan el proyecto. Debe estar relacionado directamente con la solución propuesta y su implementación, aplicable a software, hardware o infraestructura.]}

% ==========================================================
% CAPÍTULO II — METODOLOGÍA
%   Diagnóstico y área de estudio: entorno de hardware y software
%     (SPEC §2 y §5); aquí se imprime la línea/sublínea seleccionada.
%   Metodología: enfoque experimental cuantitativo con diseño factorial
%     sobre L_in, L_out, concurrencia B, chunked prefill y tasa de
%     llegada λ; experimentos E1–E5 (SPEC §8).
%   Técnicas e instrumentos: cliente con streaming, nvidia-smi/DCGM y
%     métricas Prometheus de vLLM.
%   Requerimientos RF1–RF7 y RNF1–RNF5, diseño del banco de pruebas y
%     configuración base de vLLM (SPEC §7, §11 y §12).
% ==========================================================
\chapter{Metodología}
\label{chap:metodologia}

\section{Diagnóstico de la Situación Actual}
\label{sec:cap2:diagnostico}

El método de este trabajo no se diseñó en abstracto, sino sobre una infraestructura concreta cuyas características imponen condiciones a la medición. Por esa razón, el capítulo parte de un diagnóstico de la plataforma disponible antes de describir el diseño experimental. Las limitaciones identificadas aquí reaparecen en las secciones siguientes como el porqué de buena parte de las decisiones metodológicas.

La plataforma experimental es el clúster de computación de alto rendimiento de CEDIA (\texttt{hpc.cedia.edu.ec}), gestionado con el planificador de trabajos Slurm~\cite{yoo2003slurm}.
Se eligió esta plataforma porque es la que ofrece acceso a la GPU NVIDIA A100 que fija el caso de estudio, y no por preferencia frente a otras infraestructuras. La Tabla~\ref{tab:cap2:plataforma} resume sus características. Los datos de rendimiento de la GPU proceden de la hoja de especificaciones del fabricante~\cite{nvidia2021a100}.

\begin{table}[htbp]
  \centering
  \caption{Características de la plataforma experimental.}
  \label{tab:cap2:plataforma}
  \small
  \begin{tabular}{|p{0.32\textwidth}|p{0.60\textwidth}|}
    \hline
    \textbf{Elemento} & \textbf{Valor} \\
    \hline
    Clúster & HPC CEDIA, planificador Slurm \\
    \hline
    Nodos de cómputo GPU & \texttt{compute-0-1} y \texttt{compute-0-2}, con 8 NVIDIA A100-SXM4-40GB cada uno \\
    \hline
    Nodo interactivo & \texttt{compute-0-0}, con 1 A100-40GB, 6 particiones MIG \texttt{a100\_2g.10gb} y 2 particiones MIG \texttt{a100\_1g.5gb} \\
    \hline
    CPU por nodo & 256 núcleos \\
    \hline
    RAM por nodo & $\approx$960~GB (nodos GPU); $\approx$1~TB (nodo interactivo) \\
    \hline
    Memoria de la GPU & 40\,960~MiB \\
    \hline
    Límite de potencia de la GPU & 400~W \\
    \hline
    Ancho de banda HBM2 & 1555~GB/s \\
    \hline
    Rendimiento BF16 con Tensor Cores & 312~TFLOPS (denso) \\
    \hline
    Sistema operativo & Ubuntu 20.04.6, glibc 2.31 \\
    \hline
    Almacenamiento de trabajo & NFS compartido (\texttt{\$HOME}) \\
    \hline
  \end{tabular}
\end{table}
\elaboradopor{José Andrés Ureta Granda}
\fuente{Inspección de los nodos del clúster HPC CEDIA; hoja de especificaciones de la A100~\cite{nvidia2021a100}}

La variante de la GPU es la A100-SXM4 de 40~GB, dato que la ficha institucional exige documentar porque la capacidad de memoria condiciona los niveles máximos de longitud de contexto y concurrencia. El sistema operativo de los nodos es Ubuntu 20.04.6, mientras que la ficha sugiere Ubuntu 22.04 o equivalente. La diferencia no es una elección del trabajo, sino una condición del clúster que no puede modificarse. Su consecuencia más importante es la versión de glibc, la 2.31, que origina el problema de compatibilidad tratado en la Sección~\ref{sec:cap2:parche}.

El nodo interactivo ofrece particiones MIG (\textit{Multi-Instance GPU}), que dividen una A100 en instancias aisladas de menor capacidad. Estas particiones quedan excluidas de la campaña de medición por dos razones. La primera es que su memoria resulta insuficiente para alojar Qwen3-8B. La segunda es que NVML no expone la potencia por instancia MIG, lo que impediría atribuir el consumo energético que exige el tercer objetivo específico.\footnote{La documentación de NVML define \texttt{nvmlDeviceGetPowerUsage} y \texttt{nvmlDeviceGetTotalEnergyConsumption} sobre la GPU completa y sus circuitos asociados, sin declarar soporte para los manejadores de instancia MIG; en cambio, consultas como la de memoria sí indican que la información por instancia se obtiene con esos manejadores. NVIDIA, \textit{NVML API Reference Guide}, disponible en \url{https://docs.nvidia.com/deploy/nvml-api/}. Consultado el 24 de octubre de 2026.}

La exclusión de MIG obliga, además, a una precaución de identificación. En un trabajo asignado a una partición MIG, \texttt{nvidia-smi} muestra la A100 de 40~GB que la contiene y no la instancia efectiva. Por ello, el dispositivo real se identifica con \texttt{nvidia-smi -L}, \texttt{torch.cuda.\allowbreak get\_device\_name()} o \texttt{torch.cuda.\allowbreak mem\_get\_info()}, y no con la salida general de \texttt{nvidia-smi}.

El planificador impone los límites recogidos en la Tabla~\ref{tab:cap2:slurm}. A ellos se suman dos restricciones de la calidad de servicio (QoS) \texttt{normal} asignada al usuario: solo se admite un trabajo en ejecución simultánea y no es posible reservar un nodo completo. La partición por defecto es \texttt{interactive}, que solo ofrece particiones MIG. En consecuencia, omitir la partición en la solicitud conduce a un recurso inservible para este experimento.

\begin{table}[htbp]
  \centering
  \caption{Límites de las particiones de Slurm disponibles para el usuario.}
  \label{tab:cap2:slurm}
  \small
  \begin{tabular}{|p{0.26\textwidth}|p{0.14\textwidth}|p{0.14\textwidth}|p{0.14\textwidth}|p{0.14\textwidth}|}
    \hline
    \textbf{Partición} & \textbf{Tiempo máx.} & \textbf{GPU máx.} & \textbf{CPU máx.} & \textbf{RAM máx.} \\
    \hline
    \texttt{interactive} (por defecto) & 1 día & --- & --- & --- \\
    \hline
    \texttt{gpu} & 2 días & 2 & 160 & 128~G \\
    \hline
    \texttt{gpu-max} & 1 día & 4 & 128 & 256~G \\
    \hline
    \texttt{gpu-dev} & 4 días & 1 & 16 & --- \\
    \hline
  \end{tabular}
\end{table}
\elaboradopor{José Andrés Ureta Granda}
\fuente{Configuración de Slurm del clúster HPC CEDIA}

Dos limitaciones adicionales afectan directamente a la validez de las mediciones de latencia. La primera es que los nodos son compartidos con otros usuarios y su carga de CPU varía durante la ejecución, como se detalla en la Sección~\ref{sec:cap2:control}. La segunda es que el driver de la GPU cambió durante la semana de preparación, según muestra la Tabla~\ref{tab:cap2:driver}. Este cambio indica que el clúster actualiza los drivers sin aviso previo, de modo que dos mediciones separadas por días pueden no ser comparables.

\begin{table}[htbp]
  \centering
  \caption{Versión del driver NVIDIA observada en los nodos durante la preparación.}
  \label{tab:cap2:driver}
  \small
  \begin{tabular}{|p{0.20\textwidth}|p{0.40\textwidth}|p{0.25\textwidth}|}
    \hline
    \textbf{Fecha} & \textbf{Nodo} & \textbf{Driver} \\
    \hline
    20 de octubre & \texttt{compute-0-1} & 535.161.08 \\
    \hline
    22 de octubre & \texttt{compute-0-2} & 575.51.03 \\
    \hline
    24 de octubre & \texttt{compute-0-1} y \texttt{compute-0-2} & 575.51.03 \\
    \hline
  \end{tabular}
\end{table}
\elaboradopor{José Andrés Ureta Granda}
\fuente{Registros \texttt{entorno.txt} de los trabajos ejecutados}

El diagnóstico deja establecidas, por tanto, las condiciones de partida: una A100 de 40~GB sin reserva exclusiva del nodo, un sistema operativo con glibc 2.31, un driver que puede cambiar entre ejecuciones y un único trabajo simultáneo. Antes de describir cómo se diseñó el experimento sobre esas condiciones, la Sección~\ref{sec:cap2:area} delimita el área de estudio.

\section{Área de Estudio}
\label{sec:cap2:area}

La Tabla~\ref{tab:area_estudio} sitúa el trabajo en su dominio académico y precisa el objeto de estudio. Su contenido se deriva de la ficha institucional aprobada.

\begin{table}[htbp]
  \centering
  \caption{Área de estudio.}
  \label{tab:area_estudio}
  \renewcommand{\arraystretch}{1.35}
  \begin{tabular}{|p{0.39\textwidth}|p{0.53\textwidth}|}
    \hline
    \textbf{Dominio:} & Sistemas de inferencia de modelos de lenguaje de gran tamaño sobre infraestructura de computación de alto rendimiento \\
    \hline
    \textbf{Línea de investigación:} & \UTILineaInvestigacion \\
    \hline
    \textbf{Sublínea de investigación:} & \UTISublineaInvestigacion \\
    \hline
    \textbf{Campo:} & Tecnologías de la información y la comunicación (TIC) \\
    \hline
    \textbf{Área:} & Caracterización y modelado del rendimiento de la inferencia de LLM en GPU \\
    \hline
    \textbf{Aspectos:} & Comportamiento de las fases prefill y decode; TTFT, TPOT, throughput, memoria, utilización y energía estimada; efecto de la cuantización, la longitud de contexto y la concurrencia \\
    \hline
    \textbf{Objeto de estudio:} & Modelo Qwen3-8B servido con vLLM sobre una GPU NVIDIA A100-SXM4-40GB del clúster HPC CEDIA \\
    \hline
    % PENDIENTE: fijar el periodo de estudio cuando se cierre la campaña definitiva.
    \textbf{Periodo de estudio:} & \textit{[Indique las fechas o el intervalo de tiempo que comprende el estudio.]} \\
    \hline
  \end{tabular}
\end{table}

Con el objeto de estudio delimitado, la Sección~\ref{sec:cap2:metodologia} describe el método con que se caracteriza.

\section{Metodología de Desarrollo o Implementación}
\label{sec:cap2:metodologia}

El trabajo sigue un método experimental cuantitativo organizado en las tres fases que establece la ficha institucional: preparación del entorno, benchmark controlado e instrumentación. Se adoptó esa estructura porque la ficha es el documento normativo del trabajo y porque separa con claridad lo que debe congelarse, lo que se varía y lo que se mide. La preparación del entorno se desarrolla en la Sección~\ref{sec:cap2:diseno}. Esta sección describe cómo se organiza la medición y cómo las restricciones del clúster condicionaron su diseño.

Las restricciones de Slurm determinaron tres decisiones de diseño. La primera es la ejecución secuencial de las configuraciones: con un único trabajo simultáneo, no es posible paralelizarlas entre GPU. La segunda es la elección de la partición \texttt{gpu-dev} en lugar de \texttt{gpu}. El 24 de octubre, \texttt{sbatch -{}-test-only} estimaba en \texttt{gpu} un inicio para el 31 de octubre pese a haber GPU libres, mientras que en \texttt{gpu-dev} la estimación era inmediata; \texttt{gpu-dev} ofrece, además, 4 días de tiempo máximo de ejecución frente a 2.

La tercera decisión es la reutilización del motor de inferencia. Arrancar vLLM tiene un costo fijo de decenas de segundos, que superó los 80~s en el peor caso observado, con el entorno virtual alojado en NFS. Reiniciar el motor para cada configuración multiplicaría ese costo por el número de configuraciones dentro de un único trabajo con tiempo limitado. Por ello, la fase en línea levanta el servidor una sola vez y recorre todas las configuraciones contra él.

La medición se automatiza con el script \texttt{bench\_pd.sh}, que organiza cada trabajo en las cinco fases de la Tabla~\ref{tab:cap2:fases}. Se automatizó todo el recorrido en un único script porque cada trabajo de Slurm debe ser autosuficiente: el investigador no interviene entre fases, y cualquier paso manual introduciría variaciones no registradas. Cada fase deja sus salidas en un directorio propio del trabajo, \texttt{resultados/}\textit{jobid}\texttt{/}, de modo que toda medición queda asociada al identificador del trabajo que la produjo.

\begin{table}[htbp]
  \centering
  \caption{Fases de ejecución de \texttt{bench\_pd.sh}.}
  \label{tab:cap2:fases}
  \small
  \begin{tabular}{|p{0.07\textwidth}|p{0.45\textwidth}|p{0.38\textwidth}|}
    \hline
    \textbf{Fase} & \textbf{Contenido} & \textbf{Salidas} \\
    \hline
    0 & Registro del nodo, la GPU, las versiones y la carga; arranque del monitor de telemetría & \texttt{entorno.txt}, \texttt{gpu\_monitor.csv} \\
    \hline
    1 & Prueba de humo con el modelo objetivo & \texttt{smoke.txt}, \texttt{smoke.err} \\
    \hline
    2 & Barrido sin servidor (\texttt{vllm bench latency}) & \texttt{lat\_in*\_out*\_bs*.json} \\
    \hline
    3 & Barrido en línea (\texttt{vllm serve} y \texttt{vllm bench serve}) & \texttt{serve\_in*\_out*\_c*.json}, \texttt{metrics\_servidor.txt} \\
    \hline
    4 & Agregación y tablas & \texttt{RESUMEN.md} \\
    \hline
  \end{tabular}
\end{table}
\elaboradopor{José Andrés Ureta Granda}
\fuente{Elaboración propia a partir de \texttt{bench\_pd.sh}}

El orden de las fases responde a un criterio de economía del tiempo de GPU. La fase 0 registra el entorno antes de medir, porque el driver y la carga del nodo pueden cambiar entre trabajos. La fase 1 verifica que el modelo carga y genera texto antes de comprometer el resto del tiempo asignado, de modo que un fallo del entorno se detecta antes de invertir horas de GPU en un barrido inválido. Las fases 2 y 3 aportan dos perspectivas complementarias de la misma medición, que la Sección~\ref{sec:cap2:tecnicas} explica.

El barrido implementado tiene carácter de barrido piloto y no reproduce todavía la matriz de la ficha. La ficha prevé este paso en la mitigación de su riesgo R6: ejecutar un piloto que cronometre cada configuración y renegociar la matriz con el tutor a partir de ese dato. La Tabla~\ref{tab:cap2:divergencias} declara todas las divergencias entre ambas matrices. El piloto explora un espacio más amplio en longitud de entrada y concurrencia para calibrar la matriz definitiva, que se acordará por escrito con el tutor antes de la campaña.

\begin{table}[htbp]
  \centering
  \caption{Divergencias entre la matriz de la ficha y el barrido piloto implementado.}
  \label{tab:cap2:divergencias}
  \small
  \begin{tabular}{|p{0.22\textwidth}|p{0.30\textwidth}|p{0.38\textwidth}|}
    \hline
    \textbf{Variable} & \textbf{Ficha aprobada} & \textbf{Barrido piloto (\texttt{bench\_pd.sh})} \\
    \hline
    Variante numérica & BF16 y AWQ INT4 & Solo BF16 \\
    \hline
    Longitud de entrada & 128, 512, 2048 y 8192 & 128, 512, 1024, 2048 y 4096 \\
    \hline
    Longitud de salida & 64 y 256 & 1 y 129 (sin servidor); 128 y 256 (en línea) \\
    \hline
    Concurrencia & 1, 4, 8 y 16 & 1, 2, 4, 8, 16, 32 y 64 \\
    \hline
    Repeticiones & 5 & 10 iteraciones (sin servidor) \\
    \hline
  \end{tabular}
\end{table}
\elaboradopor{José Andrés Ureta Granda}
\fuente{Ficha institucional ITI-TFG-B26-0017 y \texttt{bench\_pd.sh}}

Dos de esas divergencias son huecos que el piloto debe cerrar antes de fijar la matriz. El primero es el escenario de 8192 tokens, que no se mide. Es la configuración límite en memoria, según el cálculo de la Sección~\ref{sec:cap2:techos}, y por tanto la que decide si la matriz de la ficha es ejecutable en una A100 de 40~GB; puede exigir, además, desactivar los grafos CUDA con \texttt{-{}-enforce-eager}. El segundo hueco es la variante AWQ INT4~\cite{lin2024awq}, que constituye la mitad de la comparación exigida por el objetivo general y cuyo riesgo asociado, la disponibilidad y fiabilidad de un checkpoint cuantizado, sigue abierto.

Quedan otras dos decisiones pendientes que la memoria declara en lugar de omitir. La ficha prevé aleatorizar el orden de las configuraciones para reducir el efecto de la temperatura, la caché y el estado previo del sistema, y el material del piloto no documenta esa aleatorización; la campaña definitiva deberá incorporarla. Por otra parte, el riesgo R4 de la ficha advierte que el chunked prefill puede mezclar ambas fases en un mismo lote~\cite{agrawal2024sarathi}, y la decisión de desactivarlo o documentarlo sigue abierta. Con la organización de la medición descrita, la Sección~\ref{sec:cap2:tecnicas} precisa qué se mide y con qué instrumentos.

\section{Técnicas e Instrumentos de Recolección de Información}
\label{sec:cap2:tecnicas}

La técnica central del trabajo es la medición instrumentada de la inferencia. Las métricas por fase son el TTFT (\textit{time to first token}), que caracteriza el prefill, y el TPOT (\textit{time per output token}), que caracteriza el decode. Ninguna de las dos se observa directamente en un único instrumento sin supuestos, y por esa razón se triangulan con tres métodos independientes. Si los tres métodos convergen, la medición es robusta; si divergen, la magnitud de la discrepancia acota el error de cada uno y constituye un hallazgo que merece discusión.

\subsection{Triangulación de la separación prefill/decode}
\label{sec:cap2:triangulacion}

El primer método obtiene las fases por diferencia de latencias en el modo sin servidor, con \texttt{vllm bench latency}. Con una longitud de salida de un token, el tiempo medido corresponde esencialmente al prefill; con 129 tokens, se añaden 128 pasos de decode. Se eligió la salida de 129 tokens para que la diferencia contenga exactamente 128 pasos de decode. Las Ecuaciones~\ref{eq:cap2:tprefill} a~\ref{eq:cap2:tpot} formalizan el procedimiento, donde $T(n)$ es la latencia medida con $n$ tokens de salida, $L_{in}$ la longitud de entrada y $B$ el tamaño de lote.

\begin{equation}
\label{eq:cap2:tprefill}
T_{prefill} \approx T(1)
\end{equation}
\addeqtolist{Latencia de prefill por diferencia de latencias}

\begin{equation}
\label{eq:cap2:xprefill}
X_{prefill} = \frac{L_{in} \times B}{T(1)}
\end{equation}
\addeqtolist{Throughput de prefill sin servidor}

\begin{equation}
\label{eq:cap2:tpot}
\mathrm{TPOT} = \frac{T(129) - T(1)}{128}
\end{equation}
\addeqtolist{TPOT por diferencia de latencias}

La ventaja de este método es que aísla las fases sin instrumentación interna del motor. Su limitación es que $T(1)$ no es prefill puro: incluye una pequeña fracción de decode y el sobrecoste de arranque de la iteración. Por ello, la Ecuación~\ref{eq:cap2:tprefill} es una aproximación y su error debe contrastarse con los otros dos métodos.

El segundo método recoge las métricas de cliente en el modo en línea, con \texttt{vllm bench serve} contra un servidor \texttt{vllm serve}. La herramienta reporta el TTFT, el TPOT y el ITL (\textit{inter-token latency}, latencia entre tokens consecutivos) con los percentiles 50, 90 y 99. Se ejecuta con la opción \texttt{-{}-ignore-eos} para que cada solicitud genere la longitud de salida pedida aunque el modelo emita antes el token de fin de secuencia; sin ella, solicitudes de la misma configuración tendrían longitudes distintas y su TPOT no sería comparable. Este método ofrece la perspectiva del usuario, y su limitación es que el TTFT incluye el tiempo de espera en cola del servidor, además del prefill.

El tercer método lee las métricas internas del servidor en el endpoint \texttt{/metrics} de vLLM. Los histogramas \texttt{request\_\allowbreak prefill\_\allowbreak time\_\allowbreak seconds} y \texttt{request\_\allowbreak decode\_\allowbreak time\_\allowbreak seconds} separan las fases según la contabilidad interna del motor. Es la medida más directa y por eso sirve de referencia para validar los dos métodos anteriores. Su limitación es que depende de cómo el motor define el inicio y el fin de cada fase, definición que no es observable desde fuera y que puede variar entre versiones de vLLM.\footnote{Las opciones de \texttt{vllm bench serve} (\texttt{-{}-ignore-eos}, \texttt{-{}-metric-percentiles}) y los histogramas \texttt{vllm:request\_prefill\_time\_seconds} y \texttt{vllm:request\_decode\_time\_seconds} se documentan en el código fuente de vLLM~0.20.0 (\texttt{vllm/benchmarks/serve.py} y \texttt{vllm/v1/metrics/loggers.py}), disponible en \url{https://github.com/vllm-project/vllm/tree/v0.20.0}. Consultado el 24 de octubre de 2026.}

\subsection{Telemetría, energía y control de la contención}
\label{sec:cap2:control}

La telemetría de la GPU se registra con un monitor de \texttt{nvidia-smi} que muestrea cada 2~s durante toda la ejecución y escribe \texttt{gpu\_monitor.csv}. Se recurrió a \texttt{nvidia-smi} porque está disponible en los nodos y expone la memoria, la utilización, la temperatura y la potencia que exige el tercer objetivo específico. La ficha prevé DCGM para la campaña principal y Nsight Systems solo en un subconjunto representativo, por el volumen de datos y el sobrecoste del perfilado.

La energía se calculará a partir de esa telemetría con las Ecuaciones~\ref{eq:cap2:energia} y~\ref{eq:cap2:etoken}, definidas en la ficha. En ellas, $P_i$ es la potencia de la muestra $i$, $\Delta t$ el periodo de muestreo, $E_{idle}$ la energía consumida en reposo durante el mismo intervalo y $N_{tokens}$ el número de tokens generados. La energía así obtenida es una estimación basada en la telemetría interna de la GPU y no una medición eléctrica calibrada externamente. La ventana de integración y la medición de $E_{idle}$ en cada campaña quedan por fijar en el protocolo de energía.

\begin{equation}
\label{eq:cap2:energia}
E = \sum_{i=1}^{n} P_i \, \Delta t
\end{equation}
\addeqtolist{Energía estimada a partir de la telemetría}

\begin{equation}
\label{eq:cap2:etoken}
E_{token} = \frac{E_{ejecucion} - E_{idle}}{N_{tokens}}
\end{equation}
\addeqtolist{Energía estimada por token generado}

La contención del nodo es una variable de confusión, porque el TTFT y el TPOT dependen de que la CPU alimente a la GPU sin demora. Su magnitud quedó documentada durante la preparación: el 22 de octubre a las 19:34, \texttt{compute-0-2} registraba un load average de 317,33, 337,72 y 214,73 en las ventanas de 1, 5 y 15 minutos, sobre 256 núcleos. Esa carga equivale a una sobresuscripción de 1,24 veces en la ventana de un minuto y de 1,32 veces en la de cinco, frente a 0,84 veces quince minutos antes. La carga se duplicó, por tanto, mientras el trabajo estaba en ejecución, y nada garantiza que no vuelva a ocurrir.

Como el usuario no puede reservar un nodo completo, la contención no puede eliminarse y se controla de tres formas. La primera es registrar \texttt{/proc/loadavg} al inicio de cada trabajo en \texttt{entorno.txt}, para que la carga quede asociada a cada medición. La segunda es el monitor de \texttt{nvidia-smi} descrito, que permite detectar intervalos de baja utilización atribuibles a la CPU. La tercera es ejecutar repeticiones suficientes para estimar la varianza y detectar ejecuciones contaminadas, en lugar de confiar en una medición única.

El cambio de driver de la Tabla~\ref{tab:cap2:driver} se controla con el mismo principio de registro. Cada ejecución guarda la versión del driver vigente en el momento de medir en \texttt{entorno.txt}. La campaña definitiva deberá completarse dentro de una ventana sin actualizaciones del driver o, si eso no es posible, reportar el driver como covariable en el análisis. Definidos los instrumentos y sus controles, la Sección~\ref{sec:cap2:requerimientos} fija los requerimientos que debe cumplir el banco de pruebas que los integra.

\section{Requerimientos del Proyecto}
\label{sec:cap2:requerimientos}

Los requerimientos del banco de pruebas se derivan de los objetivos específicos y de las condiciones identificadas en el diagnóstico. Se formularon como requerimientos verificables porque el banco es un instrumento de medición, y un instrumento sin especificación no puede validarse. Las Tablas~\ref{tab:cap2:rf} y~\ref{tab:cap2:rnf} los recogen junto con la razón de cada uno.

\subsection{Requerimientos Funcionales}
\label{sec:cap2:rf}

\begin{table}[htbp]
  \centering
  \caption{Requerimientos funcionales del banco de pruebas.}
  \label{tab:cap2:rf}
  \small
  \begin{tabular}{|p{0.07\textwidth}|p{0.43\textwidth}|p{0.40\textwidth}|}
    \hline
    \textbf{ID} & \textbf{Requerimiento} & \textbf{Razón} \\
    \hline
    RF1 & Desplegar vLLM con una configuración parametrizada en YAML y verificar su salud mediante \texttt{/health} & Evita medir contra un servidor que aún no está listo \\
    \hline
    RF2 & Generar prompts sintéticos de longitud exacta en tokens con el tokenizador de Qwen3 & La longitud de entrada es una variable experimental y debe controlarse en tokens, no en caracteres \\
    \hline
    RF3 & Ejecutar la carga en modo sin servidor (lote fijo) y en modo en línea (llegadas Poisson, cliente asíncrono con streaming) & Sostiene los métodos 1 y 2 de la triangulación \\
    \hline
    RF4 & Registrar los tiempos por petición y por token con marcas de tiempo de alta resolución en el cliente & Permite calcular TTFT, TPOT e ITL desde la perspectiva del usuario \\
    \hline
    RF5 & Muestrear en paralelo la telemetría de la GPU y las métricas Prometheus de vLLM & Sostiene el método 3, la estimación de energía y el control de la contención \\
    \hline
    RF6 & Persistir cada corrida con sus metadatos: versiones, hardware, configuración, semilla y marca de tiempo & Hace comparables mediciones separadas por días o por cambios de driver \\
    \hline
    RF7 & Ajustar los modelos y generar automáticamente las tablas y figuras para LaTeX & Elimina la transcripción manual de resultados a la memoria \\
    \hline
  \end{tabular}
\end{table}
\elaboradopor{José Andrés Ureta Granda}
\fuente{Especificación del proyecto (\texttt{SPEC.md}, §11.1)}

\subsection{Requerimientos No Funcionales}
\label{sec:cap2:rnf}

\begin{table}[htbp]
  \centering
  \caption{Requerimientos no funcionales del banco de pruebas.}
  \label{tab:cap2:rnf}
  \small
  \begin{tabular}{|p{0.20\textwidth}|p{0.38\textwidth}|p{0.32\textwidth}|}
    \hline
    \textbf{ID} & \textbf{Requerimiento} & \textbf{Razón} \\
    \hline
    RNF1 Reproducibilidad & Versiones fijadas, semillas fijas, contenedor o lockfile y scripts idempotentes & El primer objetivo específico exige un entorno reproducible \\
    \hline
    RNF2 Aislamiento & GPU dedicada sin otros procesos durante las mediciones, verificada con \texttt{nvidia-smi} & Otro proceso en la GPU contaminaría latencia y energía; la contención de CPU, que no puede aislarse, se controla según la Sección~\ref{sec:cap2:control} \\
    \hline
    RNF3 Precisión de medición & Reloj monotónico (\texttt{time.perf\_counter\_ns}), warm-up descartado y repeticiones suficientes & Un reloj ajustable o una primera ejecución en frío sesgarían el TTFT \\
    \hline
    RNF4 Trazabilidad & Cada figura o tabla de la memoria referencia el identificador de la corrida de origen & Permite auditar cualquier cifra hasta su dato crudo \\
    \hline
    RNF5 Costo & La campaña completa cabe en el presupuesto de horas de GPU & Con un único trabajo simultáneo, el tiempo de GPU es el recurso escaso \\
    \hline
  \end{tabular}
\end{table}
\elaboradopor{José Andrés Ureta Granda}
\fuente{Especificación del proyecto (\texttt{SPEC.md}, §11.2)}

El requerimiento RNF1 es el que más condiciona el diseño, porque obliga a congelar un entorno que el clúster no ofrece listo para usar. La Sección~\ref{sec:cap2:diseno} describe cómo se construyó ese entorno.

\section{Diseño de la Solución}
\label{sec:diseno_solucion}
\label{sec:cap2:diseno}

El diseño de la solución comprende dos elementos: el entorno de software congelado sobre el que se ejecutan todas las mediciones y el parche que hace posible ese entorno en los nodos del clúster. Ambos responden al requerimiento RNF1, porque una medición solo es reproducible si lo es el entorno que la produce.

\subsection{Entorno de software congelado}
\label{sec:cap2:entorno}

La Tabla~\ref{tab:cap2:entorno} recoge los componentes del entorno y sus versiones exactas. Se congelaron la revisión del modelo y el commit de la rueda de vLLM, y no solo sus números de versión, porque un mismo número de versión puede corresponder a binarios o pesos distintos. Las dependencias completas quedan en el archivo \texttt{req-congelado.txt}, que acompaña al repositorio. El entorno se construyó como un entorno Python versionado, una de las tres vías de reconstrucción que admite la ficha junto con Docker y Apptainer.

\begin{table}[htbp]
  \centering
  \caption{Entorno de software congelado.}
  \label{tab:cap2:entorno}
  \small
  \begin{tabular}{|p{0.30\textwidth}|p{0.62\textwidth}|}
    \hline
    \textbf{Componente} & \textbf{Versión} \\
    \hline
    Python & 3.11 (módulo \texttt{python/3.11}, miniconda) \\
    \hline
    PyTorch & 2.11.0+cu129 \\
    \hline
    CUDA runtime & 12.9 \\
    \hline
    vLLM & 0.20.0+cu129 \\
    \hline
    Procedencia de la rueda & \texttt{wheels.vllm.ai}, commit \texttt{88d34c6409\allowbreak e9fb3c7b8c\allowbreak a0c04756f0\allowbreak 61d2099eb1} \\
    \hline
    Backend de atención & FlashAttention 2 \\
    \hline
    Modelo & \texttt{Qwen/Qwen3-8B}, revisión \texttt{b968826d9c\allowbreak 46dd6066d1\allowbreak 09eabc6255\allowbreak 188de91218} \\
    \hline
    Caché de modelos & \texttt{HF\_HOME=\textasciitilde/tfg-jose-ureta/hf-cache} \\
    \hline
    Dependencias completas & \texttt{req-congelado.txt} \\
    \hline
  \end{tabular}
\end{table}
\elaboradopor{José Andrés Ureta Granda}
\fuente{Entorno instalado en el clúster HPC CEDIA}

La razón de cada versión, en particular la de vLLM 0.20.0 y la del backend FlashAttention 2, se expone en la Sección~\ref{sec:cap2:justificacion}. Antes es necesario resolver un problema que impedía cargar esa versión de vLLM en los nodos.

\subsection{Parche de compatibilidad glibc}
\label{sec:cap2:parche}

La rueda \texttt{vllm-0.20.0+cu129} se distribuye con la etiqueta \texttt{manylinux\_2\_31}, que declara compatibilidad con glibc 2.31 o superior~\cite{smith2019pep600}. Sin embargo, dos de sus extensiones nativas, \texttt{\_C.abi3.so} y \texttt{\_moe\_C.abi3.so}, enlazan cuatro símbolos con la versión \texttt{GLIBC\_2.34}: \texttt{dlopen}, \texttt{dlclose}, \texttt{dlsym} y \texttt{pthread\_once}. En los nodos, con glibc 2.31, el resultado es un \texttt{ImportError} al cargar vLLM. La etiqueta de la rueda es, por tanto, incorrecta.

El diagnóstico consistió en inspeccionar con \texttt{objdump -T} todas las bibliotecas compartidas del entorno virtual. Solo cuatro declaran dependencias de glibc 2.32 o superior, y la Tabla~\ref{tab:cap2:bibliotecas} indica cuáles se cargan efectivamente en el clúster. Las dos que no se cargan quedan fuera del problema: una implementa FlashAttention 3, que solo existe para Hopper, y la otra está compilada para Python 3.12. El parche debe actuar, en consecuencia, solo sobre las dos extensiones del núcleo de vLLM.

\begin{table}[htbp]
  \centering
  \caption{Bibliotecas del entorno con dependencias de glibc 2.32 o superior.}
  \label{tab:cap2:bibliotecas}
  \small
  \begin{tabular}{|p{0.52\textwidth}|p{0.38\textwidth}|}
    \hline
    \textbf{Biblioteca} & \textbf{¿Se carga en el clúster?} \\
    \hline
    \texttt{vllm/\_C.abi3.so} & Sí: núcleo de vLLM \\
    \hline
    \texttt{vllm/\_moe\_C.abi3.so} & Sí \\
    \hline
    \texttt{vllm/vllm\_flash\_attn/\allowbreak\_vllm\_fa3\_C.abi3.so} & No: FlashAttention 3 es exclusiva de Hopper \\
    \hline
    \texttt{vllm/third\_party/deep\_gemm/\allowbreak\_C.cpython-312-*.so} & No: compilada para Python 3.12 \\
    \hline
  \end{tabular}
\end{table}
\elaboradopor{José Andrés Ureta Granda}
\fuente{Inspección con \texttt{objdump -T} del entorno virtual}

Los cuatro símbolos no corresponden a funcionalidad ausente en glibc 2.31. En glibc 2.34 las bibliotecas \texttt{libdl} y \texttt{libpthread} se consolidaron dentro de \texttt{libc}; en glibc 2.31 esos mismos símbolos existen, pero residen en \texttt{libdl.so.2} y \texttt{libpthread.so.0}.\footnote{GNU C Library, notas de la versión 2.34 (archivo \texttt{NEWS}, apartado \textit{Major new features}): toda la funcionalidad de \texttt{libpthread}, \texttt{libdl}, \texttt{libutil} y \texttt{libanl} se integró en \texttt{libc}. Disponible en \url{https://sourceware.org/git/?p=glibc.git;a=blob_plain;f=NEWS;hb=glibc-2.34}. Consultado el 24 de octubre de 2026.}
La incompatibilidad es, por tanto, un artefacto de haber compilado la rueda en un sistema más moderno, y basta con indicar al cargador dinámico dónde encontrar los símbolos.

La solución se implementó en el script \texttt{parche\_glibc.sh}, que aplica tres operaciones. Primero, elimina la versión de los cuatro símbolos con \texttt{patchelf -{}-clear-symbol-version} y añade \texttt{libdl.so.2} y \texttt{libpthread.so.0} como dependencias explícitas con \texttt{-{}-add-needed}.\footnote{patchelf, utilidad para modificar ejecutables y bibliotecas ELF; ambas opciones se describen en su documentación y su código fuente, disponibles en \url{https://github.com/NixOS/patchelf}. Consultado el 24 de octubre de 2026.}
Segundo, marca como \texttt{VER\_FLG\_WEAK} la entrada \texttt{GLIBC\_2.34} de la sección \texttt{.gnu.\allowbreak version\_r} mediante \texttt{parche\_\allowbreak glibc\_\allowbreak weak.py}, porque el paso anterior no basta: el cargador sigue rechazando la biblioteca mientras esa entrada exista sin la bandera. Tercero, verifica que \texttt{import vllm.\_C} e \texttt{import vllm.\_moe\_C} funcionen.

El script es idempotente y conserva los binarios originales como \texttt{*.so.orig}, de modo que aplicarlo varias veces no altera el resultado y el cambio es reversible. Debe reaplicarse tras cada reinstalación del entorno virtual, y por eso \texttt{bench\_pd.sh} lo invoca automáticamente si detecta que vLLM no se importa. Así, ningún trabajo de medición depende de que el parche se haya aplicado a mano.

El parche no altera el comportamiento numérico de vLLM. Solo reubica dónde busca el cargador dinámico cuatro símbolos de carga de bibliotecas y sincronización; no modifica kernels, precisión numérica ni planificación de solicitudes. Aun así, constituye una modificación del binario distribuido, y el entorno resultante no es bit a bit idéntico al que se obtendría con \texttt{pip install vllm==0.20.0} en un sistema con glibc 2.34 o superior. Por esa razón, los scripts del parche acompañan al repositorio, para que cualquier persona pueda reproducir exactamente el mismo entorno. Las alternativas que se evaluaron antes de optar por el parche se discuten en la Sección~\ref{sec:cap2:justificacion}.

\section{Justificación Técnica}
\label{sec:cap2:justificacion}

La justificación técnica parte de las tecnologías que fija el caso de estudio: el motor de inferencia vLLM, que gestiona el KV cache mediante PagedAttention y planifica la ejecución a nivel de iteración~\cite{kwon2023pagedattention}; el modelo Qwen3-8B~\cite{qwen2025qwen3}; y la GPU NVIDIA A100~\cite{nvidia2021a100}. Esas tecnologías vienen dadas por la ficha. Esta sección justifica, en cambio, las decisiones que el trabajo sí tomó dentro de ellas: las versiones, el backend de atención, la vía de compatibilidad y la factibilidad de la matriz en memoria.

\subsection{Versiones y backend de atención}
\label{sec:cap2:versiones}

La versión 0.20.0 de vLLM no se eligió por preferencia, sino por una restricción del entorno. Las versiones posteriores distribuyen su stack compilado contra CUDA 13, y el clúster operaba con el driver 535 en el momento de la instalación. La compatibilidad de versiones menores de CUDA permite ejecutar runtimes 12.x sobre drivers 525 o superiores, pero exige un driver 580 o superior para CUDA 13, de modo que no cruza de la versión mayor 12 a la 13.\footnote{NVIDIA, \textit{CUDA Toolkit Release Notes}, tabla de versiones mínimas del driver para la compatibilidad de versiones menores, disponible en \url{https://docs.nvidia.com/cuda/cuda-toolkit-release-notes/}; y \textit{CUDA Compatibility}, apartado de compatibilidad hacia adelante, disponible en \url{https://docs.nvidia.com/deploy/cuda-compatibility/forward-compatibility.html}. Consultado el 24 de octubre de 2026.} La versión 0.20.0 con CUDA runtime 12.9 es, por tanto, la más reciente ejecutable con aquel driver sin recurrir a la compatibilidad hacia adelante, el mecanismo de NVIDIA que permite ejecutar CUDA 13 sobre drivers anteriores en GPU de centro de datos mediante el paquete \texttt{cuda-compat}.
% REVISAR: justificar por qué no se usó la compatibilidad hacia adelante (cuda-compat), que NVIDIA declara compatible con la rama 535, o valorar si permite una versión más reciente de vLLM

La actualización posterior del driver a 575.51.03, recogida en la Tabla~\ref{tab:cap2:driver}, no invalida esta decisión. Ese driver sigue por debajo de 580, de modo que la compatibilidad de versiones menores tampoco permitiría ejecutar CUDA 13 con él. Mantener la versión congelada preserva, además, la comparabilidad con las ejecuciones ya realizadas.

vLLM selecciona FlashAttention-2~\cite{dao2024flashattention2} como backend de atención porque FlashAttention-3 está diseñada para explotar características de la arquitectura Hopper~\cite{shah2024flashattention3}, y la A100 pertenece a la arquitectura Ampere. FlashAttention reduce los accesos a la memoria de alto ancho de banda de la GPU durante el cálculo de la atención~\cite{dao2022flashattention}.
Se trata de una limitación de la plataforma y no de una decisión del experimento. Se declara porque afecta al rendimiento absoluto de la fase de prefill, que es la más intensiva en la atención.

Para resolver la incompatibilidad con glibc se evaluaron tres alternativas antes de optar por el parche, recogidas en la Tabla~\ref{tab:cap2:alternativas}. El contenedor oficial queda como plan B porque resuelve el problema sin modificar binarios, pero su tamaño lo hace costoso en el almacenamiento compartido. La compilación desde el código fuente se descartó por su costo en tiempo, y la regresión a una versión anterior, porque no resolvía el problema.

\begin{table}[htbp]
  \centering
  \caption{Alternativas evaluadas para resolver la incompatibilidad con glibc 2.31.}
  \label{tab:cap2:alternativas}
  \small
  \begin{tabular}{|p{0.40\textwidth}|p{0.50\textwidth}|}
    \hline
    \textbf{Alternativa} & \textbf{Motivo del descarte} \\
    \hline
    Contenedor oficial \texttt{vllm/vllm-openai} con enroot/pyxis o Apptainer & Imagen superior a 10~GB sobre NFS compartido; queda como plan B \\
    \hline
    Compilar vLLM desde el código fuente en el nodo & Varias horas de compilación por cada reinstalación \\
    \hline
    Fijar una versión anterior de vLLM & No resuelve el problema: la etiqueta incorrecta aparece en varias versiones \\
    \hline
  \end{tabular}
\end{table}
\elaboradopor{José Andrés Ureta Granda}
\fuente{Elaboración propia}

Fijadas las versiones, queda por justificar que la plataforma puede alojar la matriz experimental y qué rendimiento cabe esperar de ella como máximo. La Sección~\ref{sec:cap2:techos} lo establece mediante cálculos analíticos.

\subsection{Techos teóricos y presupuesto de memoria}
\label{sec:cap2:techos}

Los techos teóricos sirven de referencia contra la cual se contrastarán las mediciones. Se calculan con el modelo roofline, que acota el rendimiento de un cálculo por el menor de dos límites: el cómputo máximo del procesador y el ancho de banda de la memoria~\cite{williams2009roofline}. Sin esta referencia, una cifra de throughput es un número suelto; con ella, indica qué fracción del límite físico alcanza cada fase.

El punto de partida es el tamaño de los pesos. Qwen3-8B tiene aproximadamente $N = 8{,}2 \cdot 10^{9}$ parámetros~\cite{qwen2025qwen3}, y en BF16 cada parámetro ocupa 2 bytes. La Ecuación~\ref{eq:cap2:pesos} obtiene el volumen de los pesos, que equivale a 15,3~GiB.

\begin{equation}
\label{eq:cap2:pesos}
M_{pesos} = N \times 2~\mathrm{B} = 8{,}2 \cdot 10^{9} \times 2~\mathrm{B} \approx 16{,}4~\mathrm{GB}
\end{equation}
\addeqtolist{Volumen de los pesos de Qwen3-8B en BF16}

En el decode con lote 1, cada token generado exige leer todos los pesos desde la memoria HBM2 de la GPU, y el cómputo por byte leído es bajo. La fase queda entonces limitada por el ancho de banda de memoria, y su techo se obtiene con la Ecuación~\ref{eq:cap2:techo-decode}. El cálculo ignora la lectura del KV cache, que con lote 1 es pequeña frente a los pesos, de modo que el techo real es algo inferior.

\begin{equation}
\label{eq:cap2:techo-decode}
X_{decode}^{max} \approx \frac{BW_{HBM}}{M_{pesos}} = \frac{1555~\mathrm{GB/s}}{16{,}4~\mathrm{GB}} \approx 95~\mathrm{tok/s}
\end{equation}
\addeqtolist{Techo teórico del decode con lote 1}

En el prefill, en cambio, todos los tokens del prompt se procesan en paralelo y los pesos leídos se reutilizan para cada uno de ellos. La fase queda limitada por el cómputo, a razón de aproximadamente dos operaciones de coma flotante por parámetro y token, una multiplicación y una suma. La Ecuación~\ref{eq:cap2:techo-prefill} da su techo con el rendimiento BF16 denso de la A100.

\begin{equation}
\label{eq:cap2:techo-prefill}
X_{prefill}^{max} \approx \frac{P_{BF16}}{2N} = \frac{312 \cdot 10^{12}~\mathrm{FLOP/s}}{2 \times 8{,}2 \cdot 10^{9}~\mathrm{FLOP/tok}} \approx 19\,000~\mathrm{tok/s}
\end{equation}
\addeqtolist{Techo teórico del prefill}

Los dos techos difieren en dos órdenes de magnitud, y esa diferencia es la expresión cuantitativa del argumento central del trabajo: prefill y decode son fases con cuellos de botella distintos. Contrastar las mediciones contra estos techos es lo que permitirá identificar el cuello de botella dominante de cada fase, como exige el cuarto objetivo específico.

El presupuesto de memoria determina si la configuración más exigente de la matriz cabe en la GPU. Los parámetros de arquitectura se tomaron del archivo \texttt{config.json} de la revisión del modelo efectivamente descargada, y no de la documentación, porque es esa revisión la que se ejecuta. El modelo tiene 36 capas y 8 cabezas KV de dimensión 128, gracias a la atención de consultas agrupadas (GQA)~\cite{ainslie2023gqa}. La Ecuación~\ref{eq:cap2:kv-token} calcula el KV cache por token, donde el factor 2 inicial corresponde a las claves y los valores, y los 2 bytes finales, a la representación BF16.

\begin{equation}
\label{eq:cap2:kv-token}
m_{KV} = 2 \times n_{capas} \times n_{KV} \times d_{h} \times b = 2 \times 36 \times 8 \times 128 \times 2~\mathrm{B} = 147\,456~\mathrm{B} = 144{,}0~\mathrm{KiB}
\end{equation}
\addeqtolist{KV cache por token de Qwen3-8B}

La configuración más exigente de la ficha combina 8192 tokens con 16 secuencias concurrentes. La Ecuación~\ref{eq:cap2:kv-total} multiplica el KV cache por token por ese volumen, y la Ecuación~\ref{eq:cap2:mem-total} le suma los pesos. El total, 33,3~GiB, queda dentro de los aproximadamente 36~GiB utilizables de la A100 de 40~GB, pero con un margen estrecho.

\begin{equation}
\label{eq:cap2:kv-total}
M_{KV} = 8192 \times 16 \times 147\,456~\mathrm{B} \approx 18{,}0~\mathrm{GiB}
\end{equation}
\addeqtolist{KV cache de la configuración más exigente}

\begin{equation}
\label{eq:cap2:mem-total}
M_{total} = M_{pesos} + M_{KV} \approx 15{,}3~\mathrm{GiB} + 18{,}0~\mathrm{GiB} = 33{,}3~\mathrm{GiB}
\end{equation}
\addeqtolist{Presupuesto de memoria de la configuración más exigente}

El cálculo indica que la matriz de la ficha es factible en principio, pero no lo garantiza: el margen restante debe alojar también las activaciones y los grafos CUDA. Por esa razón, el escenario de 8192 tokens con concurrencia 16 en BF16 se identificó en la Sección~\ref{sec:cap2:metodologia} como el hueco prioritario del piloto. Con el método, los instrumentos, el entorno y las cotas analíticas establecidos, el Capítulo~\ref{chap:implementacion} describe la implementación del banco de pruebas y la ejecución de la campaña.

% ==========================================================
% CAPÍTULO III — IMPLEMENTACIÓN, PRUEBAS Y RESULTADOS
%   Implementación del banco de pruebas (bench/ y analysis/).
%   Campaña experimental E1–E5: prefill aislado, decode aislado,
%     extremo a extremo, interferencia y capacidad.
%   Ajuste y validación hold-out de los modelos de TTFT, TPOT y E2E
%     frente a los criterios de aceptación CA1–CA6 y a las cotas
%     roofline (SPEC §8–10).
%   Cronograma y presupuesto de horas-GPU (SPEC §13).
% ==========================================================
\chapter{Implementación, Pruebas y Resultados}
\label{chap:implementacion}

\section{Descripción de la Solución}
\label{sec:cap3:descripcion}

La solución implementada es un banco de pruebas que mide por separado las fases prefill y decode de Qwen3-8B servido con vLLM en una A100-SXM4-40GB del clúster HPC CEDIA. Se compone del entorno de software congelado y del parche de compatibilidad descritos en la Sección~\ref{sec:cap2:diseno}, del script de medición \texttt{bench\_pd.sh}, cuyas fases recoge la Tabla~\ref{tab:cap2:fases}, y del script de campaña \texttt{a100\_qwen3\_8b.sh}, que la Sección~\ref{sec:cap3:rediseno} presenta como respuesta a los problemas observados en la primera ejecución.

Este capítulo informa de la primera ejecución completa sobre la A100 con el modelo objetivo, el trabajo de Slurm 30732. Ese trabajo cubrió las fases 0 a 2 de \texttt{bench\_pd.sh}: registro del entorno, prueba de humo y barrido sin servidor. Sus resultados son, por tanto, resultados piloto: validan el banco y el método de diferencia de latencias, pero no constituyen todavía la campaña experimental.
% PENDIENTE: describir el cliente tfgbench (RF3-RF6) y los experimentos E1 y E2 cuando se complete la campaña.

\section{Proceso de Implementación}
\label{sec:cap3:proceso}

\subsection{Ejecución del trabajo 30732}
\label{sec:cap3:trabajo}

El trabajo 30732 se ejecutó el 24 de octubre de 2026 sobre una A100 completa, sin particiones MIG. La Tabla~\ref{tab:cap3:trabajo} recoge el entorno registrado por la fase 0 en \texttt{entorno.txt}. La carga del nodo al inicio, 37,20 en la ventana de un minuto sobre 256 núcleos, es baja frente a la documentada en la Sección~\ref{sec:cap2:control}.

\begin{table}[htbp]
  \centering
  \caption{Entorno registrado en el trabajo 30732.}
  \label{tab:cap3:trabajo}
  \small
  \begin{tabular}{|p{0.32\textwidth}|p{0.60\textwidth}|}
    \hline
    \textbf{Elemento} & \textbf{Valor} \\
    \hline
    Nodo & \texttt{compute-0-1} \\
    \hline
    Inicio & 24 de octubre de 2026, 11:07:24 (UTC$-$5) \\
    \hline
    Carga al inicio (1, 5 y 15 min) & 37,20; 36,72; 40,91 \\
    \hline
    GPU & NVIDIA A100-SXM4-40GB, 40\,960~MiB, límite de 400~W, 1410~MHz máx. \\
    \hline
    Driver & 535.161.08 \\
    \hline
    Software & vLLM 0.20.0, PyTorch 2.11.0+cu129, CUDA 12.9 \\
    \hline
    Modelo & \texttt{Qwen/Qwen3-8B}, BF16 \\
    \hline
    Opciones del motor & \texttt{max\_model\_len}=8192, \texttt{gpu\_memory\_utilization}=0,90, caché de prefijos desactivada, chunked prefill activo con 8192 tokens por lote, grafos CUDA y \texttt{torch.compile} activos \\
    \hline
  \end{tabular}
\end{table}
\elaboradopor{José Andrés Ureta Granda}
\fuente{\texttt{entorno.txt} y registros del trabajo 30732}

Cada configuración del barrido sin servidor se midió con \texttt{vllm bench latency}, que crea un motor nuevo en cada invocación, ejecuta tres iteraciones de calentamiento y mide diez. La Tabla~\ref{tab:cap3:arranques} desglosa la duración de cada ejecución. La columna de inicialización corresponde al tiempo que vLLM declara para perfilar la memoria, crear el KV cache y calentar el modelo, del cual la mayor parte es compilación con \texttt{torch.compile}.

\begin{table}[htbp]
  \centering
  \caption{Duración de las ejecuciones del trabajo 30732.}
  \label{tab:cap3:arranques}
  \small
  \begin{tabular}{|p{0.24\textwidth}|p{0.12\textwidth}|p{0.18\textwidth}|p{0.14\textwidth}|p{0.13\textwidth}|}
    \hline
    \textbf{Ejecución} & \textbf{Inicio} & \textbf{Inicialización del motor (s)} & \textbf{Duración total (min)} & \textbf{Medición (s)} \\
    \hline
    Prueba de humo ($L_{in}$ máx. 2048) & 11:32:35 & 632,8 & 15,9 & --- \\
    \hline
    $L_{in}=128$, $L_{out}=1$ & 12:11:47 & 1452,4 & 74,6 & 0,17 \\
    \hline
    $L_{in}=128$, $L_{out}=129$ & 13:49:08 & 436,4 & 28,7 & 16,95 \\
    \hline
    $L_{in}=512$, $L_{out}=1$ & 14:47:50 & 248,6 & 43,1 & 0,37 \\
    \hline
    $L_{in}=512$, $L_{out}=129$ & 15:55:36 & --- & Incompleta & --- \\
    \hline
  \end{tabular}
\end{table}
\elaboradopor{José Andrés Ureta Granda}
\fuente{Registros \texttt{smoke.txt} y \texttt{lat\_in*\_out*\_bs1.log} del trabajo 30732}

La tabla muestra una desproporción que condiciona el resto de la implementación. Las tres configuraciones completadas sumaron menos de 18~s de medición efectiva en un trabajo de casi cinco horas. La primera compilación de \texttt{torch.compile} tardó 1416~s y guardó su resultado en la caché \texttt{\textasciitilde/.cache/\allowbreak vllm/\allowbreak torch\_compile\_cache}, lo que explica que las inicializaciones siguientes bajaran a 436~s y 249~s. Aun así, cada ejecución siguió costando entre 29 y 75 minutos.

La inicialización del motor no explica toda esa duración. En la primera ejecución transcurrieron 38 minutos entre la lectura de la configuración, a las 12:12:25, y la creación del motor, a las 12:50:35, y entre el final de una ejecución y el primer registro de la siguiente pasaron 23 y 30 minutos. La causa de estas esperas no quedó determinada. Los pesos residen en un sistema de archivos NFS4, según el propio registro de vLLM, y su lectura tardó 57~s en la primera ejecución y 10~s en las siguientes, de modo que la lectura de los pesos tampoco basta para explicarlas.

La quinta configuración no llegó a medirse. La telemetría se detiene a las 16:06:11, cuando el motor de esa ejecución aún se estaba inicializando, y ningún registro indica un error.
% REVISAR: confirmar con sacct -j 30732 si el trabajo terminó por el límite de tiempo.

\subsection{Rediseño de la ejecución}
\label{sec:cap3:rediseno}

El trabajo 30732 demostró que, en esta plataforma, el costo dominante no es la medición sino el arranque del motor. Con un arranque de entre 29 y 75 minutos, un barrido que crea un motor por configuración no cabe en un trabajo de Slurm, por largo que sea su límite de tiempo. La decisión de la Sección~\ref{sec:cap2:metodologia} de levantar el servidor una sola vez y recorrer todas las configuraciones contra él pasa así de ser una optimización a ser una condición de viabilidad.

El script \texttt{a100\_qwen3\_8b.sh} materializa esa decisión. Solicita una A100 completa en la partición \texttt{gpu-dev}, levanta un único \texttt{vllm serve} con las mismas opciones de la Tabla~\ref{tab:cap3:trabajo}, espera hasta 45 minutos a que el endpoint \texttt{/health} responda y ejecuta después los experimentos con el cliente \texttt{tfgbench} en el mismo nodo, contra \texttt{127.0.0.1}. Al ejecutar el cliente y el servidor en el mismo nodo, el TTFT medido no incluye latencia de red. Tras cada experimento, el script guarda las métricas acumuladas de \texttt{/metrics}, que sostienen el tercer método de la triangulación de la Sección~\ref{sec:cap2:triangulacion}. Además, reduce el periodo de muestreo de la telemetría de 2~s a 1~s. La Tabla~\ref{tab:cap3:e1e2} resume los dos experimentos que el script ejecuta.

\begin{table}[htbp]
  \centering
  \caption{Experimentos ejecutados por \texttt{a100\_qwen3\_8b.sh}.}
  \label{tab:cap3:e1e2}
  \small
  \begin{tabular}{|p{0.22\textwidth}|p{0.34\textwidth}|p{0.34\textwidth}|}
    \hline
    \textbf{Parámetro} & \textbf{E1: prefill} & \textbf{E2: concurrencia} \\
    \hline
    Longitud de entrada & 128, 256, 512, 1024, 2048, 4096 y 7680 & 1024 \\
    \hline
    Longitud de salida & 128 & 256 \\
    \hline
    Concurrencia & 1 & 1, 2, 4, 8, 16, 32 y 64 \\
    \hline
    Solicitudes medidas por punto & 30 (más 3 de calentamiento) & 256 (más 4 de calentamiento) \\
    \hline
    Métricas principales & TTFT & TPOT, ITL y throughput \\
    \hline
  \end{tabular}
\end{table}
\elaboradopor{José Andrés Ureta Granda}
\fuente{\texttt{configs/e1\_prefill.toml} y \texttt{configs/e2\_concurrencia.toml}}

En ambos experimentos los prompts son secuencias aleatorias de identificadores de token de longitud exacta, generadas con semilla 0 y sin prefijos compartidos, y las solicitudes se envían con temperatura 0 e \texttt{ignore\_eos}. La longitud máxima de E1 es 7680 y no 8192, porque la suma de entrada y salida no puede superar el \texttt{max\_model\_len} de 8192. En E2, con 64 secuencias de 1280 tokens, el KV cache ocupa unos 11,3~GiB y cabe en la memoria disponible sin expulsiones, según el presupuesto de la Sección~\ref{sec:cap3:memoria}.

El primer lanzamiento de este script, el trabajo 31157 del 27 de octubre, no llegó a medir. El servidor no respondió en los 45 minutos de espera, su registro quedó vacío y la GPU no llegó a ocupar memoria.
% PENDIENTE: actualizar con el resultado del diagnóstico de los trabajos 31157 y 31299 y con la ejecución de E1 y E2.

\section{Pruebas y Validación}
\label{sec:cap3:pruebas}

\subsection{Prueba de humo y verificación del dispositivo}
\label{sec:cap3:humo}

La prueba de humo cargó Qwen3-8B y generó 64 tokens de texto coherente en español a partir de un prompt sobre la inferencia de modelos de lenguaje. Su tiempo total, 13,77~s, incluye la primera ejecución del modelo y no se usa como medida de rendimiento. La prueba cumple su función: confirma que el entorno parcheado de la Sección~\ref{sec:cap2:parche} carga el modelo y genera texto antes de comprometer el resto del trabajo.

El dispositivo se verificó con \texttt{nvidia-smi -L}, que identifica una A100-SXM4-40GB completa y no una instancia MIG, conforme a la precaución de la Sección~\ref{sec:cap2:diagnostico}. El backend de atención seleccionado fue FlashAttention-2, como anticipa la Sección~\ref{sec:cap2:versiones}.

\subsection{Estabilidad de las mediciones}
\label{sec:cap3:estabilidad}

La Tabla~\ref{tab:cap3:estabilidad} resume la dispersión de las diez iteraciones medidas en cada configuración. El coeficiente de variación no supera el 0,6~\%, y el rango completo de cada configuración es inferior a 1~ms. Con lote 1 y sin servidor, la medición es, por tanto, muy repetible dentro de un mismo trabajo.

\begin{table}[htbp]
  \centering
  \caption{Dispersión de la latencia en las diez iteraciones medidas.}
  \label{tab:cap3:estabilidad}
  \small
  \begin{tabular}{|p{0.22\textwidth}|p{0.12\textwidth}|p{0.14\textwidth}|p{0.10\textwidth}|p{0.11\textwidth}|p{0.11\textwidth}|}
    \hline
    \textbf{Configuración} & \textbf{Media (ms)} & \textbf{Desv. típica (ms)} & \textbf{CV (\%)} & \textbf{Mín. (ms)} & \textbf{Máx. (ms)} \\
    \hline
    $L_{in}=128$, $L_{out}=1$ & 17,087 & 0,099 & 0,58 & 16,943 & 17,205 \\
    \hline
    $L_{in}=512$, $L_{out}=1$ & 37,207 & 0,138 & 0,37 & 36,984 & 37,396 \\
    \hline
    $L_{in}=128$, $L_{out}=129$ & 1695,221 & 0,306 & 0,02 & 1694,844 & 1695,809 \\
    \hline
  \end{tabular}
\end{table}
\elaboradopor{José Andrés Ureta Granda}
\fuente{Archivos \texttt{lat\_in*\_out*\_bs1.json} del trabajo 30732}

Las iteraciones de calentamiento confirman la necesidad de descartarlas. La primera de cada configuración tardó 25,2~s, 14,9~s y 33,5~s, respectivamente, entre 9 y 1500 veces más que la media de las iteraciones medidas, porque incluye la captura de grafos y la carga de kernels pendientes. Incluirla en la media habría invalidado la medición, como prevé el requerimiento RNF3.

La baja dispersión no garantiza, sin embargo, la reproducibilidad entre trabajos. Las diez iteraciones de cada configuración se ejecutaron en menos de 20~s, dentro de un mismo motor y con una carga del nodo estable. La variación entre trabajos, nodos y versiones del driver, que la Sección~\ref{sec:cap2:control} identifica como la principal amenaza, solo podrá estimarse con repeticiones en trabajos distintos.

\subsection{Contraste del presupuesto de memoria}
\label{sec:cap3:memoria}

Los registros de vLLM permiten contrastar el presupuesto analítico de la Sección~\ref{sec:cap2:techos}. Los pesos ocuparon 15,27~GiB, frente a los 15,3~GiB de la Ecuación~\ref{eq:cap2:pesos}. Con \texttt{gpu\_memory\_utilization} de 0,90, vLLM reservó 16,32~GiB para el KV cache, equivalentes a 118\,832 tokens, y 2,44~GiB para los grafos CUDA. El cociente entre ambas cifras del KV cache es de 147\,456~bytes por token, que coincide exactamente con la Ecuación~\ref{eq:cap2:kv-token}.

El contraste revela que la configuración más exigente de la ficha no cabe en la GPU con estas opciones. Con 8192 tokens por solicitud, el KV cache disponible admite 14,51 secuencias simultáneas, según el propio registro de vLLM, y no las 16 que exige la ficha. En las ejecuciones posteriores, la reserva subió a 16,85~GiB y 122\,672 tokens, que admiten 14,97 secuencias, cifra que sigue por debajo de 16. El presupuesto de la Ecuación~\ref{eq:cap2:mem-total} era correcto en sus términos, pero omitía los 2,44~GiB de los grafos CUDA y el 10~\% de la memoria que vLLM deja sin reservar.

Este resultado cierra el primero de los huecos señalados en la Sección~\ref{sec:cap2:metodologia}. La combinación de 8192 tokens y concurrencia 16 solo sería ejecutable aumentando \texttt{gpu\_memory\_utilization}, desactivando los grafos CUDA con \texttt{-{}-enforce-eager} o aceptando que el planificador ponga solicitudes en espera. Cualquiera de las tres opciones altera lo que se mide, y por esa razón E1 limita la entrada a 7680 tokens con una sola solicitud, y E2 limita la entrada a 1024 tokens.
% PENDIENTE: acordar con el tutor la sustitución del nivel 8192 x 16 de la ficha.

\section{Resultados}
\label{sec:cap3:resultados}

Los resultados de esta sección proceden exclusivamente del barrido sin servidor del trabajo 30732, con lote $B=1$. Se calculan con el método de diferencia de latencias de la Sección~\ref{sec:cap2:triangulacion} y se contrastan con los techos teóricos de la Sección~\ref{sec:cap2:techos}.

\subsection{Fase de prefill}
\label{sec:cap3:prefill}

La Tabla~\ref{tab:cap3:prefill} aplica las Ecuaciones~\ref{eq:cap2:tprefill} y~\ref{eq:cap2:xprefill} a las dos longitudes de entrada medidas. La columna de MFU (\textit{model FLOPs utilization}) expresa el throughput de prefill como fracción del techo de cómputo~\cite{chowdhery2023palm}, con la Ecuación~\ref{eq:cap3:mfu}, donde $P_{BF16}$ es el rendimiento BF16 denso de la A100 y $N$ el número de parámetros.

\begin{equation}
\label{eq:cap3:mfu}
\mathrm{MFU} = \frac{2N \times X_{prefill}}{P_{BF16}}
\end{equation}
\addeqtolist{Utilización de cómputo del prefill}

\begin{table}[htbp]
  \centering
  \caption{Latencia y throughput de prefill con lote 1.}
  \label{tab:cap3:prefill}
  \small
  \begin{tabular}{|p{0.16\textwidth}|p{0.22\textwidth}|p{0.26\textwidth}|p{0.18\textwidth}|}
    \hline
    \textbf{$L_{in}$ (tokens)} & \textbf{$T_{prefill}$ (ms)} & \textbf{$X_{prefill}$ (tok/s)} & \textbf{MFU (\%)} \\
    \hline
    128 & 17,09 & 7\,491 & 39,4 \\
    \hline
    512 & 37,21 & 13\,761 & 72,3 \\
    \hline
  \end{tabular}
\end{table}
\elaboradopor{José Andrés Ureta Granda}
\fuente{Elaboración propia a partir del trabajo 30732}

El throughput de prefill casi se duplica al cuadruplicar la entrada, mientras que la latencia solo se multiplica por 2,2. Ese comportamiento indica que una parte de $T_{prefill}$ no depende de la longitud de entrada. Si se supone una relación lineal, $T_{prefill} = a + b \, L_{in}$, los dos puntos medidos determinan la pendiente de la Ecuación~\ref{eq:cap3:pendiente} y una ordenada en el origen $a \approx 10{,}4$~ms.

\begin{equation}
\label{eq:cap3:pendiente}
b = \frac{T_{prefill}(512) - T_{prefill}(128)}{512 - 128} = \frac{37{,}21 - 17{,}09}{384}~\mathrm{ms} \approx 0{,}0524~\mathrm{ms/token}
\end{equation}
\addeqtolist{Costo marginal del prefill por token de entrada}

El inverso de esa pendiente equivale a 19\,085~tok/s, prácticamente igual al techo de 19\,000~tok/s de la Ecuación~\ref{eq:cap2:techo-prefill}. Sin embargo, esta coincidencia no permite afirmar que el prefill alcance el 100~\% del techo de cómputo en el margen. Lo más probable es que el modelo lineal sea demasiado simple. Con 128 tokens, las multiplicaciones de matrices son pequeñas y aprovechan peor los Tensor Cores que con 512, de modo que la eficiencia crece con $L_{in}$ y la pendiente entre dos puntos queda subestimada. Además, la Ecuación~\ref{eq:cap2:techo-prefill} ignora las operaciones de la atención, que crecen con el cuadrado de la longitud.

Con dos puntos no es posible separar el costo fijo de la variación de la eficiencia. Por esa razón, el experimento E1 mide siete longitudes entre 128 y 7680 tokens. El dato robusto de este piloto es la cota inferior: con 512 tokens, el prefill ya alcanza el 72~\% del techo de cómputo, lo que confirma que es una fase limitada por el cómputo.

\subsection{Fase de decode}
\label{sec:cap3:decode}

La Ecuación~\ref{eq:cap2:tpot} aplicada a la entrada de 128 tokens da un TPOT de 13,11~ms, equivalente a 76,3~tok/s por secuencia. Ese valor representa el 80~\% del techo de 95~tok/s de la Ecuación~\ref{eq:cap2:techo-decode}. La Ecuación~\ref{eq:cap3:mbu} expresa la misma relación como MBU (\textit{model bandwidth utilization}), la fracción del ancho de banda de la memoria que se aprovecha al leer los pesos en cada paso~\cite{agarwal2023llminference}.

\begin{equation}
\label{eq:cap3:mbu}
\mathrm{MBU} = \frac{M_{pesos} / \mathrm{TPOT}}{BW_{HBM}} = \frac{16{,}4~\mathrm{GB} / 13{,}11~\mathrm{ms}}{1555~\mathrm{GB/s}} = \frac{1251~\mathrm{GB/s}}{1555~\mathrm{GB/s}} \approx 80~\%
\end{equation}
\addeqtolist{Utilización del ancho de banda en el decode}

Un MBU del 80~\% indica que el decode con lote 1 está limitado por el ancho de banda de la memoria y opera cerca de su techo. El 20~\% restante corresponde a la lectura del KV cache, que la Ecuación~\ref{eq:cap2:techo-decode} ignora, y al sobrecosto por paso del motor, que no se puede separar con estos datos. La desviación típica de 0,31~ms sobre 1695~ms de la Tabla~\ref{tab:cap3:estabilidad} implica que el TPOT se conoce con una incertidumbre inferior a 0,01~ms.

El TPOT de 512 tokens de entrada no se obtuvo porque la ejecución con $L_{out}=129$ no se completó. Por tanto, el piloto no permite observar cómo crece el TPOT con el contexto, que es el efecto esperado de la lectura del KV cache.

\subsection{Contraste entre fases}
\label{sec:cap3:contraste}

Los dos resultados anteriores cuantifican la asimetría que motiva el trabajo. Procesar los 128 tokens del prompt costó 17,09~ms, mientras que generar un único token costó 13,11~ms. En la práctica, el prefill procesó 128 tokens en un tiempo apenas un 30~\% mayor que el que el decode necesita para generar uno solo. La Tabla~\ref{tab:cap3:contraste} resume ambas fases frente a sus techos.

\begin{table}[htbp]
  \centering
  \caption{Comparación de las fases frente a sus techos teóricos ($L_{in}=128$, $B=1$).}
  \label{tab:cap3:contraste}
  \small
  \begin{tabular}{|p{0.13\textwidth}|p{0.19\textwidth}|p{0.17\textwidth}|p{0.15\textwidth}|p{0.16\textwidth}|}
    \hline
    \textbf{Fase} & \textbf{Recurso limitante} & \textbf{Medido (tok/s)} & \textbf{Techo (tok/s)} & \textbf{Utilización} \\
    \hline
    Prefill & Cómputo & 7\,491 ($L_{in}=128$) \newline 13\,761 ($L_{in}=512$) & 19\,000 & MFU 39~\% \newline MFU 72~\% \\
    \hline
    Decode & Ancho de banda & 76,3 & 95 & MBU 80~\% \\
    \hline
  \end{tabular}
\end{table}
\elaboradopor{José Andrés Ureta Granda}
\fuente{Elaboración propia a partir del trabajo 30732 y de la Sección~\ref{sec:cap2:techos}}

Las dos fases difieren en dos órdenes de magnitud en throughput por secuencia, como anticipaba el análisis de la Sección~\ref{sec:cap2:techos}. Con lote 1, el decode aprovecha una fracción mayor de su recurso limitante que el prefill de entradas cortas. Sin embargo, su throughput absoluto es 98 veces menor con 128 tokens y 180 veces menor con 512.

\subsection{Telemetría de la GPU}
\label{sec:cap3:telemetria}

La GPU ocupó como máximo 36\,799~MiB, el 90~\% de su memoria, conforme al valor de \texttt{gpu\_memory\_utilization}. En reposo, con la memoria vacía, la potencia media fue de 55,2~W. La potencia máxima, 268~W, se registró durante el calentamiento y la medición de la configuración con 129 tokens de salida, que es la única que mantuvo la GPU ocupada durante más de unos segundos.

La telemetría del trabajo 30732 no permite estimar la energía por token con la Ecuación~\ref{eq:cap2:etoken}. Con un muestreo cada 2~s, las mediciones de prefill, que duraron menos de medio segundo, no contienen ninguna muestra propia, y la de decode contiene unas ocho. Esta limitación motivó la reducción del periodo a 1~s en \texttt{a100\_qwen3\_8b.sh}. En E2, cuyos puntos duran varios minutos, ese periodo bastará para integrar la energía, pero en E1 seguirá siendo insuficiente.
% PENDIENTE: valorar el contador de energía de NVML (nvmlDeviceGetTotalEnergyConsumption) para E1.

\subsection{Alcance de los resultados piloto}
\label{sec:cap3:alcance}

Los resultados de esta sección validan el método de diferencia de latencias y los techos teóricos, pero tienen un alcance limitado. Proceden de un único trabajo en un único nodo, con lote 1, dos longitudes de entrada y una de salida. Además, solo aplican el primero de los tres métodos de triangulación, porque la fase en línea no llegó a ejecutarse. Las métricas de cliente y las de \texttt{/metrics}, que permitirían contrastarlos, corresponden a los experimentos E1 y E2.

\section{Cronograma y Presupuesto}
\label{sec:cap3:cronograma}
% PENDIENTE: completar cuando se cierre la campaña.
\textit{[Presentar la planificación de las fases de ejecución del proyecto y el detalle de los costos asociados a hardware, licencias, servicios u otros recursos.]}

% ==========================================================
% CAPÍTULO IV — CONCLUSIONES Y RECOMENDACIONES
%   Conclusiones por objetivo específico (OE1–OE4) y recomendaciones.
%   Trabajos futuros: cuantización (FP8/AWQ), multi-GPU, desagregación
%     prefill/decode, otros motores y otros modelos.
% ==========================================================
\chapter{Conclusiones y Recomendaciones}

\section{Conclusiones}
\textit{[Redactar conclusiones basadas en los resultados obtenidos y en el cumplimiento de los objetivos planteados.]}

\section{Recomendaciones}
\textit{[Proponer mejoras, optimizaciones o ampliaciones del proyecto, fundamentadas en el análisis realizado.]}

\section{Trabajos Futuros}
\textit{[Plantear líneas de desarrollo o investigación que puedan derivarse del proyecto implementado.]}

% ==========================================================
% REFERENCIAS (capítulo NO numerado, en índice, IEEE desde referencias.bib)
% ==========================================================
\clearpage
\phantomsection
\addcontentsline{toc}{chapter}{REFERENCIAS} % para que aparezca en el índice

\noindent\textit{[Incluya todas las referencias bibliográficas utilizadas en el documento. Deben estar citadas según normas IEEE y gestionadas mediante un gestor bibliográfico como Mendeley o Zotero.]}

\bibliographystyle{IEEEtran} % formato IEEE (recomendado)
\bibliography{referencias}    % usa referencias.bib
% ==========================================================
% GLOSARIO (capítulo NO numerado)
%   Términos y siglas del trabajo: prefill, decode, TTFT, TPOT, ITL,
%   caché KV, PagedAttention, continuous batching, chunked prefill,
%   goodput, MFU, MBU, roofline, etc.
% ==========================================================

\clearpage
\chapter*{GLOSARIO}
\addcontentsline{toc}{chapter}{GLOSARIO}

\noindent\textit{[Defina los términos técnicos relevantes utilizados en el documento.]}

\begin{description}
  \item[Término o sigla 1] \textit{[Definición técnica breve y clara.]}
  \item[Término o sigla 2] \textit{[Definición técnica breve y clara.]}
\end{description}
% ==========================================================
% ANEXOS (ANEXO A, ANEXO B, ... + Índice de Anexos)
%   Incluir el enlace al repositorio con el código, los datos crudos
%   y el entorno versionado (criterio de aceptación CA6).
% ==========================================================
\UTIAppendixMatter

\Anexo{Título del Anexo A}

\textit{[Incluya una descripción breve del contenido del anexo. Ejemplo: documentación técnica, diagramas, configuraciones, evidencias o instrumentos aplicados.]}

\begin{itemize}
  \item \textit{[Documentación técnica.]}
  \item \textit{[Manuales de usuario y administración.]}
  \item \textit{[Evidencias de pruebas.]}
  \item \textit{[Autorizaciones y validaciones de la entidad receptora.]}
  \item \textit{[Repositorio de versionamiento (enlace).]}
  \item \textit{[Formatos de entrevistas y encuestas aplicadas.]}
\end{itemize}

% \Anexo{Título del Anexo B}
% Contenido...

\end{document}
